\chapter[Effectful Mealy Morphisms and Finite-Path Feedback]
{Effectful Mealy Morphisms, Fibred Centrality, and Finite-Path Feedback}
\label{chap:effectful-mealy-centrality-feedback}

\section{Orientation}
\label{sec:effectful-orientation}

The Mealy morphisms developed in Chapter~2 are deterministic.  For every
admissible input--state pair
\[
    (a,x)\in A_1\times_{t_A,A_0,p}P,
\]
the structure map selects one updated state and one output-comparison arrow.
The effectful replacement keeps these pieces of data together and places the
resulting pair inside one computation:
\[
    \widehat\Phi:
    P\circ A
    \longrightarrow
    \mathsf K_{A_0\times B_0}(B\circ P).
\]
Thus an outcome of processing
\[
    a:u\to v=p(x)
\]
is a joint state--output value
\[
    (x',\beta),
    \qquad
    p(x')=u,
    \qquad
    \beta:q(x')\to q(x),
\]
and different computational branches may contain different updated states and
different output arrows.

The effect must preserve the external interface of the carrier.  It therefore
acts in the slice over the interface object rather than merely on the apex of
a span.  We fix a finitely complete category \(\mathcal E\), with chosen
pullbacks, and a fibred strong monad
\[
    \mathsf K
\]
on the codomain fibration
\[
    \operatorname{cod}:\mathcal E^{\to}\longrightarrow\mathcal E.
\]
For every object \(Z\), this gives a strong monad
\[
    \mathsf K_Z:
    \mathcal E/Z\longrightarrow\mathcal E/Z,
\]
and pullback reindexing commutes coherently with these slice monads.

A strong monad carries two canonical orders for combining computations.  They
need not agree.  Their discrepancy is the obstruction to middle-four
interchange for effectful maps of spans.  Fibrewise central computations
commute with all computations in one slice.  Span composition, however,
reindexes computations to larger interface products.  The correct class for
the span calculus is therefore the class of computations whose every
reindexing is fibrewise central.  We call these computations
\emph{fibred-central}.  They form the largest reindexing-stable part of the
Kleisli calculus on which the two execution orders agree.

The chapter has five principal constructions.

\begin{enumerate}
    \item It constructs the indexed Kleisli calculus induced by
    \(\mathsf K\), including dependent-sum transport, its compatibility with
    strength, and the resulting projection formula.

    \item It constructs the bicategory of spans with fibred-central effectful
    2-cells.

    \item It constructs the bicategory of fibred-central effectful Mealy
    morphisms and effectful simulations.

    \item It constructs a quotient of unpointed state presentations by
    zigzags of strict state retractions, proves that this quotient is not
    forced to collapse through the empty presentation, and establishes its
    universal retract-invariant state-hiding property.

    \item Under the feedback-base conditions of Chapter~9, it lifts homwise
    sums, the coproduct tensor, finite path objects, and finite hidden-path
    trace to the effectful setting.
\end{enumerate}

The unrestricted connected-component quotient by all strict state
homomorphisms exists, but is unusably coarse in the unpointed setting.  As
soon as the base has a pullback-stable initial object, the empty state
presentation maps strictly into every other presentation, so all
presentations with fixed source and target become connected.  The
feedback-base conditions of Chapter~9 include pullback stability of the
empty coproduct and hence make the initial object strict and pullback-stable.
The replacement used below therefore identifies only presentations connected
by strict state retractions.

This construction is related to, but distinct from, the pointed
bisimilarity quotient used to present uniform feedback in effectful-machine
semantics
\cite{BonchiDiLavoreRomanEffectfulMealyMachines}.  That construction is
pointed: a machine carries an initial state, and machine homomorphisms
preserve that point.  A Mealy morphism in the present manuscript carries an
unpointed state object.  The universal property established below is
therefore a universal property for retract-invariant unpointed state hiding,
not for pointed uniform feedback.

Throughout this chapter, equations involving span composition, pullback,
dependent sum, and reindexing are read under the canonical chosen pullback
associators, unitors, Beck--Chevalley isomorphisms, Frobenius isomorphisms,
and their coherence equations.  As in the preceding chapters, all
generalized-element formulas denote the corresponding morphism equations out
of the appropriate pullback objects.

\section{Effects over interfaces}
\label{sec:effects-over-interfaces}

\subsection{The codomain bifibration}
\label{subsec:codomain-bifibration-effects}

\begin{definition}[Fibred strong monad]
\label{def:fibred-strong-monad}
A \textbf{fibred strong monad} on the codomain fibration consists of a monad
\[
    (\mathsf K,\eta,\mu)
\]
on \(\mathcal E^{\to}\) lying over the identity functor of \(\mathcal E\),
such that:
\begin{enumerate}
    \item \(\mathsf K\) preserves cartesian arrows;

    \item \(\eta\) and \(\mu\) are vertical natural transformations;

    \item every slice restriction
    \[
        \mathsf K_Z:\mathcal E/Z\to\mathcal E/Z
    \]
    carries a left strength
    \[
        \operatorname{st}^{\ell,Z}_{P,Q}:
        P\times_Z\mathsf K_ZQ
        \longrightarrow
        \mathsf K_Z(P\times_ZQ);
    \]

    \item the strengths satisfy the strong-monad equations and commute
    coherently with pullback reindexing.
\end{enumerate}
\end{definition}

Thus the monad, its strength, and all reindexing coherence are given once on
the codomain fibration.  They are not selected independently for each
interface.

For a morphism
\[
    h:W\longrightarrow Z,
\]
write
\[
    h^*:\mathcal E/Z\longrightarrow\mathcal E/W
\]
for pullback.  Postcomposition with \(h\) defines a left adjoint
\[
    \Sigma_h:\mathcal E/W\longrightarrow\mathcal E/Z.
\]
Thus the codomain fibration is also an opfibration, and there is an
adjunction
\[
    \Sigma_h\dashv h^*.
\]
For \(Q\in\mathcal E/W\), write
\[
    \nu^h_Q:
    Q\longrightarrow h^*\Sigma_hQ
\]
for the unit of this adjunction.

Fibredness supplies, for every \(Q\in\mathcal E/Z\), an invertible comparison
\[
    \varphi_{h,Q}:
    h^*\mathsf K_ZQ
    \xrightarrow{\cong}
    \mathsf K_Wh^*Q,
\]
natural in \(h\) and \(Q\), pseudofunctorial in \(h\), and compatible with
\(\eta\), \(\mu\), and the strengths.

For a Kleisli arrow
\[
    f:P\longrightarrow\mathsf K_ZQ,
\]
define its reindexing by
\[
    h^{*\mathsf K}(f)
    :=
    \varphi_{h,Q}\circ h^*f:
    h^*P\longrightarrow\mathsf K_Wh^*Q.
\]

\begin{definition}[Dependent-sum transport of effects]
\label{def:dependent-sum-effect-transport}
For
\[
    h:W\to Z
\]
and \(Q\in\mathcal E/W\), define
\[
    \kappa_{h,Q}:
    \Sigma_h\mathsf K_WQ
    \longrightarrow
    \mathsf K_Z\Sigma_hQ
\]
to be the mate, under \(\Sigma_h\dashv h^*\), of the composite
\[
    \mathsf K_WQ
    \xrightarrow{\mathsf K_W(\nu^h_Q)}
    \mathsf K_Wh^*\Sigma_hQ
    \xrightarrow{\varphi_{h,\Sigma_hQ}^{-1}}
    h^*\mathsf K_Z\Sigma_hQ.
\]
\end{definition}

The map \(\kappa_{h,Q}\) moves a computation through the operation that
forgets the intermediate interface \(W\to Z\).  It is forced by fibredness
and the adjunction \(\Sigma_h\dashv h^*\).

\begin{proposition}[Functoriality of dependent-sum transport]
\label{prop:dependent-sum-effect-functoriality}
Let
\[
    W\xrightarrow{h}Z\xrightarrow{g}V
\]
be morphisms in \(\mathcal E\).  Under the canonical identifications
\(\Sigma_{gh}\cong\Sigma_g\Sigma_h\), one has
\[
    \kappa_{1_Z,Q}=1_{\mathsf K_ZQ}
\]
and
\[
    \kappa_{gh,Q}
    =
    \kappa_{g,\Sigma_hQ}
    \circ
    \Sigma_g(\kappa_{h,Q}).
\]
Moreover,
\begin{align}
    \kappa_{h,Q}\circ\Sigma_h(\eta^W_Q)
    &=
    \eta^Z_{\Sigma_hQ},
    \label{eq:kappa-unit}
    \\
    \kappa_{h,Q}\circ\Sigma_h(\mu^W_Q)
    &=
    \mu^Z_{\Sigma_hQ}
    \circ
    \mathsf K_Z(\kappa_{h,Q})
    \circ
    \kappa_{h,\mathsf K_WQ}.
    \label{eq:kappa-multiplication}
\end{align}
\end{proposition}

\begin{proof}
Each equation is proved by transposition along the appropriate
dependent-sum--pullback adjunction.

For the identity equation, the transpose of \(\kappa_{1_Z,Q}\) is
\[
    \mathsf K_ZQ
    \xrightarrow{\mathsf K_Z(\nu^{1_Z}_Q)}
    \mathsf K_ZQ
    \xrightarrow{\varphi^{-1}_{1_Z,Q}}
    \mathsf K_ZQ.
\]
Both maps are identities under the chosen identity cleavage, so the mate is
the identity.

For composition, the unit of the composite adjunction satisfies
\[
    \nu^{gh}_Q
    =
    h^*(\nu^g_{\Sigma_hQ})\circ\nu^h_Q
\]
under the canonical identification
\[
    (gh)^*\Sigma_{gh}
    \cong
    h^*g^*\Sigma_g\Sigma_h.
\]
Substitution into the defining mate of \(\kappa_{gh,Q}\), followed by
pseudofunctoriality of \(\varphi\), yields the mate of
\[
    \kappa_{g,\Sigma_hQ}
    \circ
    \Sigma_g(\kappa_{h,Q}).
\]

For the unit equation, the transpose of the left-hand side is
\[
    Q
    \xrightarrow{\eta^W_Q}
    \mathsf K_WQ
    \xrightarrow{\mathsf K_W(\nu^h_Q)}
    \mathsf K_Wh^*\Sigma_hQ
    \xrightarrow{\varphi^{-1}}
    h^*\mathsf K_Z\Sigma_hQ.
\]
Naturality of \(\eta\) and compatibility of \(\varphi\) with the monad unit
identify this with
\[
    Q
    \xrightarrow{\nu^h_Q}
    h^*\Sigma_hQ
    \xrightarrow{h^*(\eta^Z)}
    h^*\mathsf K_Z\Sigma_hQ,
\]
which is the transpose of the right-hand side.

For the multiplication equation, transpose both sides.  The transpose of
the left-hand side is
\[
\begin{aligned}
    \mathsf K_W\mathsf K_WQ
    &\xrightarrow{\mu^W_Q}
    \mathsf K_WQ
    \xrightarrow{\mathsf K_W(\nu^h_Q)}
    \mathsf K_Wh^*\Sigma_hQ
    \xrightarrow{\varphi^{-1}}
    h^*\mathsf K_Z\Sigma_hQ.
\end{aligned}
\]
The transpose of the right-hand side expands to the same map by naturality
of \(\mu\), compatibility of \(\varphi\) with multiplication, and the monad
associativity equation.
\end{proof}

\begin{corollary}[Dependent-sum Kleisli functor]
\label{cor:dependent-sum-kleisli-functor}
Every map \(h:W\to Z\) induces a functor
\[
    \Sigma_h^{\mathsf K}:
    \operatorname{Kl}(\mathsf K_W)
    \longrightarrow
    \operatorname{Kl}(\mathsf K_Z)
\]
which agrees with \(\Sigma_h\) on objects and sends
\[
    f:P\longrightarrow\mathsf K_WQ
\]
to
\[
    \Sigma_h^{\mathsf K}(f)
    :=
    \Sigma_hP
    \xrightarrow{\Sigma_hf}
    \Sigma_h\mathsf K_WQ
    \xrightarrow{\kappa_{h,Q}}
    \mathsf K_Z\Sigma_hQ.
\]
\end{corollary}

\begin{proof}
Equation~\eqref{eq:kappa-unit} gives preservation of Kleisli identities, and
Equation~\eqref{eq:kappa-multiplication} gives preservation of Kleisli
composition.
\end{proof}

\begin{proposition}[Effectful Beck--Chevalley]
\label{prop:effectful-beck-chevalley}
Consider a pullback square
\[
\begin{tikzcd}
    W'
        \arrow[r,"k'"]
        \arrow[d,"h'"']
    &
    W
        \arrow[d,"h"]
    \\
    Z'
        \arrow[r,"k"']
    &
    Z.
\end{tikzcd}
\]
There is a canonical natural isomorphism of Kleisli functors
\[
    k^{*\mathsf K}\Sigma_h^{\mathsf K}
    \xRightarrow{\cong}
    \Sigma_{h'}^{\mathsf K}k'^{*\mathsf K}.
\]
\end{proposition}

\begin{proof}
Ordinary Beck--Chevalley for the codomain bifibration gives
\[
    k^*\Sigma_h
    \xRightarrow{\cong}
    \Sigma_{h'}k'^*.
\]
To verify that this comparison intertwines the transported Kleisli arrows,
transpose both possible composites along
\(\Sigma_{h'}\dashv h'^*\).  The two transposes coincide by the defining mate
equation for \(\kappa\), pseudofunctoriality of \(\varphi\), and the equality
of the two pullback composites in the given square.
\end{proof}

\subsection{Strength and the two Fubini maps}
\label{subsec:slice-strength-fubini}

Each slice \(\mathcal E/Z\) is cartesian.  Write
\[
    \operatorname{st}^{\ell,Z}_{P,Q}:
    P\times_Z\mathsf K_ZQ
    \longrightarrow
    \mathsf K_Z(P\times_ZQ)
\]
for the left strength.  The symmetry of the slice product defines the right
strength
\[
    \operatorname{st}^{r,Z}_{P,Q}:
    \mathsf K_ZP\times_ZQ
    \longrightarrow
    \mathsf K_Z(P\times_ZQ).
\]

\begin{definition}[Slice Fubini maps]
\label{def:slice-fubini-maps}
For \(P,Q\in\mathcal E/Z\), define
\[
    d^{\ell,Z}_{P,Q},
    d^{r,Z}_{P,Q}:
    \mathsf K_ZP\times_Z\mathsf K_ZQ
    \longrightarrow
    \mathsf K_Z(P\times_ZQ)
\]
by
\begin{align*}
    d^{\ell,Z}_{P,Q}
    &:=
    \mu^Z_{P\times_ZQ}
    \circ
    \mathsf K_Z(\operatorname{st}^{r,Z}_{P,Q})
    \circ
    \operatorname{st}^{\ell,Z}_{\mathsf K_ZP,Q},
    \\
    d^{r,Z}_{P,Q}
    &:=
    \mu^Z_{P\times_ZQ}
    \circ
    \mathsf K_Z(\operatorname{st}^{\ell,Z}_{P,Q})
    \circ
    \operatorname{st}^{r,Z}_{P,\mathsf K_ZQ}.
\end{align*}
\end{definition}

The map \(d^r\) executes the \(P\)-computation before the
\(Q\)-computation.  The map \(d^\ell\) executes the \(Q\)-computation before
the \(P\)-computation.  The superscripts record the outer strength in the
displayed construction rather than the execution order.  Both maps return
their values in the same ordered product \(P\times_ZQ\).

The two maps are related by the symmetry:
\[
    d^{\ell,Z}_{P,Q}
    =
    \mathsf K_Z(\tau_{Q,P})
    \circ
    d^{r,Z}_{Q,P}
    \circ
    \tau_{\mathsf K_ZP,\mathsf K_ZQ},
\]
where \(\tau\) denotes the appropriate symmetry in the slice.

\begin{proposition}[Reindexing of Fubini maps]
\label{prop:fubini-reindexing}
For every \(h:W\to Z\), the comparison \(\varphi\) identifies
\[
    h^*(d^{\ell,Z}_{P,Q})
\]
with
\[
    d^{\ell,W}_{h^*P,h^*Q},
\]
and similarly for \(d^r\).
\end{proposition}

\begin{proof}
Pullback preserves the cartesian products in the slices.  Apply the
reindexing compatibility of the left strength, the induced right strength,
and the monad multiplication to the two composites of
Definition~\ref{def:slice-fubini-maps}.
\end{proof}

For \(h:W\to Z\), \(P\in\mathcal E/W\), and \(R\in\mathcal E/Z\), let
\[
    \operatorname{Fr}^{r}_{h;P,R}:
    \Sigma_h(P\times_Wh^*R)
    \xrightarrow{\cong}
    \Sigma_hP\times_ZR
\]
denote the right Frobenius isomorphism.  Let
\[
    \operatorname{Fr}^{\ell}_{h;R,P}:
    \Sigma_h(h^*R\times_WP)
    \xrightarrow{\cong}
    R\times_Z\Sigma_hP
\]
denote its conjugate by the slice symmetries.

\begin{lemma}[Strength compatibility of dependent-sum transport]
\label{lem:dependent-sum-strength-compatibility}
Let \(h:W\to Z\).

For \(Q\in\mathcal E/W\) and \(S\in\mathcal E/Z\), one has
\[
\begin{aligned}
&
\mathsf K_Z(\operatorname{Fr}^{r}_{h;Q,S})
\circ
\kappa_{h,Q\times_Wh^*S}
\circ
\Sigma_h
\bigl(
    \operatorname{st}^{r,W}_{Q,h^*S}
\bigr)
\\
&\qquad=
\operatorname{st}^{r,Z}_{\Sigma_hQ,S}
\circ
\bigl(
    \kappa_{h,Q}\times_Z1_S
\bigr)
\circ
\operatorname{Fr}^{r}_{h;\mathsf K_WQ,S}.
\end{aligned}
\tag{\(\mathrm{St}^r_h\)}
\]

For \(R\in\mathcal E/Z\) and \(Q\in\mathcal E/W\), one has
\[
\begin{aligned}
&
\mathsf K_Z(\operatorname{Fr}^{\ell}_{h;R,Q})
\circ
\kappa_{h,h^*R\times_WQ}
\circ
\Sigma_h
\bigl(
    \operatorname{st}^{\ell,W}_{h^*R,Q}
\bigr)
\\
&\qquad=
\operatorname{st}^{\ell,Z}_{R,\Sigma_hQ}
\circ
\bigl(
    1_R\times_Z\kappa_{h,Q}
\bigr)
\circ
\operatorname{Fr}^{\ell}_{h;R,\mathsf K_WQ}.
\end{aligned}
\tag{\(\mathrm{St}^\ell_h\)}
\]
\end{lemma}

\begin{proof}
Transpose the first equation along \(\Sigma_h\dashv h^*\).  Under the
Frobenius and pullback-product identifications, both transposes have domain
\[
    \mathsf K_WQ\times_Wh^*S
\]
and codomain
\[
    \mathsf K_Wh^*(\Sigma_hQ\times_ZS).
\]
The transpose of the left-hand side is obtained by applying the right
strength and then the mate defining \(\kappa\).  The transpose of the
right-hand side is obtained by first applying the mate defining
\(\kappa_{h,Q}\) and then the reindexed right strength.  These are equal by
the compatibility of \(\varphi\) with strength and naturality of the unit
\(\nu^h\).

The left-strength equation is proved identically.  Equivalently, it follows
from the right-strength equation by conjugating with the slice symmetries.
\end{proof}

\begin{lemma}[Projection formula for effectful dependent sums]
\label{lem:effectful-projection-formula}
Let
\[
    f:P\to\mathsf K_WQ,
    \qquad
    g:R\to\mathsf K_ZS.
\]
For \(\varepsilon\in\{\ell,r\}\), put
\[
    f\otimes^\varepsilon h^{*\mathsf K}g
    :=
    d^{\varepsilon,W}_{Q,h^*S}
    \circ
    \bigl(
        f\times_Wh^{*\mathsf K}g
    \bigr).
\]
Then
\[
\begin{aligned}
&
\mathsf K_Z(\operatorname{Fr}^{r}_{h;Q,S})
\circ
\Sigma_h^{\mathsf K}
\bigl(
    f\otimes^\varepsilon h^{*\mathsf K}g
\bigr)
\circ
(\operatorname{Fr}^{r}_{h;P,R})^{-1}
\\
&\qquad=
d^{\varepsilon,Z}_{\Sigma_hQ,S}
\circ
\bigl(
    \Sigma_h^{\mathsf K}f\times_Zg
\bigr).
\end{aligned}
\tag{\(\mathrm{PF}_\varepsilon\)}
\]
\end{lemma}

\begin{proof}
Expand \(d^\varepsilon\) into strengths and multiplication.  Move the first
strength through \(\Sigma_h\) by
Lemma~\ref{lem:dependent-sum-strength-compatibility}.  Move the second
strength through \(\Sigma_h\) by the same lemma, using the left or right
version as appropriate.  Finally move multiplication through \(\Sigma_h\)
by Equation~\eqref{eq:kappa-multiplication}.

Equivalently, transpose both sides along \(\Sigma_h\dashv h^*\).  After
applying \(\varphi\), both transposes have domain
\[
    P\times_Wh^*R
\]
and codomain
\[
    \mathsf K_Wh^*(\Sigma_hQ\times_ZS).
\]
The transpose of the left-hand side is
\[
\begin{aligned}
P\times_Wh^*R
&\xrightarrow{
    f\times h^{*\mathsf K}g
}
\mathsf K_WQ\times_W\mathsf K_Wh^*S
\\
&\xrightarrow{
    d^{\varepsilon,W}
}
\mathsf K_W(Q\times_Wh^*S)
\\
&\xrightarrow{
    \mathsf K_W(\nu^h_Q\times 1)
}
\mathsf K_W
\bigl(
    h^*\Sigma_hQ\times_Wh^*S
\bigr)
\\
&\xrightarrow{\cong}
\mathsf K_Wh^*(\Sigma_hQ\times_ZS).
\end{aligned}
\]
The transpose of the right-hand side is the same composite by
Proposition~\ref{prop:fubini-reindexing},
Lemma~\ref{lem:dependent-sum-strength-compatibility}, and naturality of the
Frobenius isomorphism.
\end{proof}

\subsection{Pure parameters in the slice Kleisli calculus}
\label{subsec:pure-parameter-kleisli-calculus}

For a Kleisli arrow
\[
    f:P\longrightarrow\mathsf K_ZQ
\]
and an object \(R\in\mathcal E/Z\), define
\[
    f\rtimes_ZR
    :=
    \operatorname{st}^{r,Z}_{Q,R}
    \circ
    (f\times_Z1_R):
    P\times_ZR
    \longrightarrow
    \mathsf K_Z(Q\times_ZR)
\]
and
\[
    R\ltimes_Zf
    :=
    \operatorname{st}^{\ell,Z}_{R,Q}
    \circ
    (1_R\times_Zf):
    R\times_ZP
    \longrightarrow
    \mathsf K_Z(R\times_ZQ).
\]

\begin{proposition}[Pure-parameter Kleisli functors]
\label{prop:pure-parameter-kleisli-functors}
For every \(R\in\mathcal E/Z\), the assignments
\[
    P\longmapsto P\times_ZR,
    \qquad
    f\longmapsto f\rtimes_ZR
\]
and
\[
    P\longmapsto R\times_ZP,
    \qquad
    f\longmapsto R\ltimes_Zf
\]
define endofunctors of
\[
    \operatorname{Kl}(\mathsf K_Z).
\]

For every map \(u:R\to R'\) in \(\mathcal E/Z\), the ordinary product maps
\[
    1_Q\times_Zu,
    \qquad
    u\times_Z1_Q
\]
define Kleisli-natural transformations between the corresponding
pure-parameter functors.
\end{proposition}

\begin{proof}
For the right pure-parameter functor, the strength-unit equation gives
\[
    \operatorname{st}^{r,Z}_{P,R}
    \circ
    (\eta^Z_P\times_Z1_R)
    =
    \eta^Z_{P\times_ZR},
\]
so Kleisli identities are preserved.

Let
\[
    f:P\to\mathsf K_ZQ,
    \qquad
    g:Q\to\mathsf K_ZS.
\]
Then
\[
    g\bullet f
    =
    \mu^Z_S
    \circ
    \mathsf K_Z(g)
    \circ
    f.
\]
Expanding
\[
    (g\bullet f)\rtimes_ZR
\]
and applying naturality of the right strength together with its compatibility
with \(\mu\) gives
\[
    (g\rtimes_ZR)
    \bullet
    (f\rtimes_ZR).
\]
Thus right pure-parameter extension preserves Kleisli composition.

The proof for left pure-parameter extension is identical using the left
strength.  Naturality in \(u\) is precisely naturality of the corresponding
strength.
\end{proof}

\begin{proposition}[Coherence of pure parameters]
\label{prop:pure-parameter-coherence}
The canonical associators, unitors, and symmetries of the slice products give
coherent Kleisli-natural isomorphisms
\[
    (f\rtimes_ZR)\rtimes_ZS
    \cong
    f\rtimes_Z(R\times_ZS),
\]
\[
    R\ltimes_Z(S\ltimes_Zf)
    \cong
    (R\times_ZS)\ltimes_Zf,
\]
and
\[
    R\ltimes_Z(f\rtimes_ZS)
    \cong
    (R\ltimes_Zf)\rtimes_ZS.
\]
\end{proposition}

\begin{proof}
The first two comparisons are the associativity equations for the right and
left strengths.  The mixed comparison follows from the standard compatibility
of the left and right strengths induced by the symmetry of the cartesian
slice.  Since only one factor is effectful, no exchange of two computations
occurs.  The coherence diagrams reduce to cartesian coherence and the
strong-functor coherence of \(\mathsf K_Z\).
\end{proof}

\section{Central computations}
\label{sec:central-computations}

\subsection{Fibrewise and fibred centrality}
\label{subsec:fibred-central-kleisli-arrows}

\begin{definition}[Fibrewise central Kleisli arrow]
\label{def:fibrewise-central-kleisli-arrow}
Let \(P,Q\in\mathcal E/Z\).  A Kleisli arrow
\[
    f:P\longrightarrow\mathsf K_ZQ
\]
is \textbf{fibrewise central} if, for every Kleisli arrow
\[
    g:R\longrightarrow\mathsf K_ZS,
\]
one has
\[
    d^{\ell,Z}_{Q,S}
    \circ
    (f\times_Zg)
    =
    d^{r,Z}_{Q,S}
    \circ
    (f\times_Zg).
\]
\end{definition}

By the symmetry relation between \(d^\ell\) and \(d^r\), a fibrewise central
arrow also commutes with an arbitrary computation when it occurs in the
second coordinate.

If \(f\) and \(g\) are fibrewise central, write
\[
    f\otimes_Zg
    :=
    d^{\ell,Z}\circ(f\times_Zg)
    =
    d^{r,Z}\circ(f\times_Zg).
\]

\begin{proposition}[Local closure of central arrows]
\label{prop:local-central-arrow-closure}
In every slice \(\mathcal E/Z\):
\begin{enumerate}
    \item every pure Kleisli arrow
    \[
        \eta^Z_Q\circ u:P\to\mathsf K_ZQ
    \]
    is fibrewise central;

    \item fibrewise central arrows are closed under Kleisli composition;

    \item if \(f\) and \(g\) are fibrewise central, then
    \(f\otimes_Zg\) is fibrewise central;

    \item the fibrewise central arrows form a wide symmetric monoidal
    subcategory
    \[
        \operatorname{Kl}^{\mathrm{loccen}}(\mathsf K_Z)
        \subseteq
        \operatorname{Kl}(\mathsf K_Z).
    \]
\end{enumerate}
\end{proposition}

\begin{proof}
For a pure arrow, the strength-unit equations reduce either Fubini order to
the image under \(\mathsf K_Z\) of the corresponding ordinary product map.
Hence the two orders agree.

Let
\[
    f:P\to\mathsf K_ZQ,
    \qquad
    f':Q\to\mathsf K_ZQ'
\]
be fibrewise central, and let
\[
    g:R\to\mathsf K_ZS
\]
be arbitrary.  Expand the two Fubini composites for
\[
    f'\mathbin{\bullet}f
    =
    \mu\circ\mathsf K(f')\circ f.
\]
In the first expansion, commute \(f\) past \(g\), and then commute \(f'\)
past \(g\).  The strong-monad associativity equations identify the resulting
composite with the opposite Fubini expansion.  Thus
\(f'\mathbin{\bullet}f\) is fibrewise central.

Now let \(f\) and \(g\) be fibrewise central and let
\[
    h:T\to\mathsf K_ZU
\]
be arbitrary.  To commute \(f\otimes_Zg\) past \(h\), first commute \(g\)
past \(h\), and then commute \(f\) past \(h\).  Associativity and symmetry of
the cartesian product, together with the associativity equations for the
Fubini maps, identify the resulting composite with the opposite execution
order.  Hence \(f\otimes_Zg\) is central.

The monoidal unit is the terminal object of the slice with its pure identity
arrow.  The associator, unitors, and symmetry are pure, and the preceding
closure results make the tensor bifunctorial on the central arrows.
\end{proof}

Fibrewise centrality alone is not sufficient for span composition, because
horizontal composition reindexes cells to larger interface products.

\begin{definition}[Fibred-central Kleisli arrow]
\label{def:fibred-central-kleisli-arrow}
A Kleisli arrow
\[
    f:P\longrightarrow\mathsf K_ZQ
\]
is \textbf{fibred-central} if, for every morphism
\[
    h:W\to Z,
\]
the reindexed Kleisli arrow
\[
    h^{*\mathsf K}f:
    h^*P\longrightarrow\mathsf K_Wh^*Q
\]
is fibrewise central.

From this point onward, the unqualified word \textbf{central} means
fibred-central.
\end{definition}

Write
\[
    \operatorname{Kl}^{\mathrm{cen}}(\mathsf K_Z)
\]
for the wide subcategory of \(\operatorname{Kl}(\mathsf K_Z)\) consisting of
central arrows.

\begin{proposition}[Closure of fibred-central arrows]
\label{prop:fibred-central-arrow-closure}
Central arrows satisfy the following closure properties.
\begin{enumerate}
    \item Every pure Kleisli arrow is central.

    \item Central arrows are closed under Kleisli composition.

    \item If \(f\) and \(g\) are central, then their common tensor
    \(f\otimes_Zg\) is central.

    \item Pullback reindexing preserves central arrows.

    \item Dependent-sum transport preserves central arrows: if
    \[
        f:P\to\mathsf K_WQ
    \]
    is central, then
    \[
        \Sigma_h^{\mathsf K}f:
        \Sigma_hP\to\mathsf K_Z\Sigma_hQ
    \]
    is central.
\end{enumerate}
\end{proposition}

\begin{proof}
The first three assertions follow by applying
Proposition~\ref{prop:local-central-arrow-closure} after every reindexing.

For pullback stability, let \(f\) be central over \(Z\), let \(h:W\to Z\),
and let \(k:V\to W\).  Pseudofunctoriality gives
\[
    k^{*\mathsf K}h^{*\mathsf K}f
    \cong
    (hk)^{*\mathsf K}f,
\]
which is fibrewise central by
Definition~\ref{def:fibred-central-kleisli-arrow}.

For dependent sums, let \(k:Z'\to Z\), and form the pullback square
\[
\begin{tikzcd}
    W'
        \arrow[r,"k'"]
        \arrow[d,"h'"']
    &
    W
        \arrow[d,"h"]
    \\
    Z'
        \arrow[r,"k"']
    &
    Z.
\end{tikzcd}
\]
By Proposition~\ref{prop:effectful-beck-chevalley},
\[
    k^{*\mathsf K}\Sigma_h^{\mathsf K}f
    \cong
    \Sigma_{h'}^{\mathsf K}k'^{*\mathsf K}f.
\]
The arrow \(k'^{*\mathsf K}f\) is fibrewise central.  Let
\[
    g:R\to\mathsf K_{Z'}S
\]
be arbitrary.  Apply
Lemma~\ref{lem:effectful-projection-formula} to
\(k'^{*\mathsf K}f\) and \(g\).  The two execution orders for
\(\Sigma_{h'}^{\mathsf K}k'^{*\mathsf K}f\) and \(g\) are the dependent-sum
images of the two execution orders for \(k'^{*\mathsf K}f\) and
\(h'^{*\mathsf K}g\).  Those two orders agree by fibrewise centrality.
Therefore every reindexing of \(\Sigma_h^{\mathsf K}f\) is fibrewise
central.
\end{proof}

\begin{corollary}[Indexed premonoidal centre]
\label{cor:indexed-premonoidal-centre}
The assignment
\[
    Z\longmapsto
    \operatorname{Kl}^{\mathrm{cen}}(\mathsf K_Z)
\]
is closed under pullback reindexing and dependent-sum transport.  Each fibre
is symmetric monoidal under the slice product, and the pure Kleisli functor
\[
    \mathcal E/Z
    \longrightarrow
    \operatorname{Kl}(\mathsf K_Z)
\]
corestricts to it.
\end{corollary}

\section{Effectful maps of spans}
\label{sec:effectful-maps-of-spans}

For a span
\[
    P=
    \left(
        X\xleftarrow{p}P\xrightarrow{q}Y
    \right),
\]
write
\[
    \overline P
    :=
    \langle p,q\rangle:
    P\longrightarrow X\times Y
\]
for the corresponding object of \(\mathcal E/(X\times Y)\).

\begin{definition}[Effectful span cell]
\label{def:effectful-span-cell}
Let \(P,Q:X\rightsquigarrow Y\) be parallel spans.  A
\textbf{\(\mathsf K\)-effectful span cell}
\[
    \widehat f:
    P\Rrightarrow_{\mathsf K}Q
\]
is a morphism
\[
    \widehat f:
    \overline P
    \longrightarrow
    \mathsf K_{X\times Y}(\overline Q)
\]
in \(\mathcal E/(X\times Y)\).

It is \textbf{central} when it is central in the sense of
Definition~\ref{def:fibred-central-kleisli-arrow}.
\end{definition}

Vertical composition is Kleisli composition:
\[
    \widehat g\bullet\widehat f
    :=
    \mu\circ\mathsf K(\widehat g)\circ\widehat f.
\]
The vertical identity on \(P\) is
\[
    \eta_{\overline P}:
    \overline P\longrightarrow
    \mathsf K_{X\times Y}(\overline P).
\]

\subsection{Ordered horizontal composition}
\label{subsec:ordered-horizontal-effectful-spans}

Let
\[
    P,P':X\rightsquigarrow Y,
    \qquad
    Q,Q':Y\rightsquigarrow Z,
\]
and let
\[
    \widehat f:P\Rrightarrow_{\mathsf K}P',
    \qquad
    \widehat g:Q\Rrightarrow_{\mathsf K}Q'.
\]
Put
\[
    W:=X\times Y\times Z,
\]
with projections
\[
    \pi_{XY}:W\to X\times Y,
    \qquad
    \pi_{YZ}:W\to Y\times Z,
    \qquad
    \pi_{XZ}:W\to X\times Z.
\]
Define
\[
    \widetilde P:=\pi_{XY}^*\overline P,
    \qquad
    \widetilde Q:=\pi_{YZ}^*\overline Q,
\]
and similarly for \(P'\) and \(Q'\).  Write
\[
    \widetilde f
    :=
    \pi_{XY}^{*\mathsf K}\widehat f,
    \qquad
    \widetilde g
    :=
    \pi_{YZ}^{*\mathsf K}\widehat g.
\]

There are canonical isomorphisms
\[
    \Sigma_{\pi_{XZ}}
    (\widetilde P\times_W\widetilde Q)
    \cong
    \overline{Q\circ P}
\]
and
\[
    \Sigma_{\pi_{XZ}}
    (\widetilde P'\times_W\widetilde Q')
    \cong
    \overline{Q'\circ P'}.
\]

\begin{definition}[Ordered horizontal composites]
\label{def:ordered-horizontal-effectful-composites}
The two ordered horizontal composites
\[
    \widehat g\ast_\ell\widehat f,
    \qquad
    \widehat g\ast_r\widehat f:
    Q\circ P
    \Rrightarrow_{\mathsf K}
    Q'\circ P'
\]
are, under the canonical identifications above,
\[
\begin{aligned}
    \widehat g\ast_\varepsilon\widehat f
    :=
    \kappa_{\pi_{XZ}}
    \circ
    \Sigma_{\pi_{XZ}}
    \bigl(
        d^{\varepsilon,W}
        \circ
        (\widetilde f\times_W\widetilde g)
    \bigr),
    \qquad
    \varepsilon\in\{\ell,r\}.
\end{aligned}
\]
\end{definition}

Thus both computations are first pulled to the complete typed interface
\(X\times Y\times Z\).  They are combined using one of the two Fubini maps,
and the intermediate \(Y\)-interface is then forgotten by
\(\Sigma_{\pi_{XZ}}^{\mathsf K}\).

\subsection{Functorial pure whiskering}
\label{subsec:functorial-pure-whiskering}

Let
\[
    Q:Y\rightsquigarrow Z
\]
be an ordinary span and let
\[
    \widehat f:
    P\Rrightarrow_{\mathsf K}P'
\]
be an effectful cell between spans \(X\rightsquigarrow Y\).

Define the pure left whiskering
\[
    Q\blacktriangleright\widehat f
\]
by pulling \(\widehat f\) and \(Q\) to \(X\times Y\times Z\), forming
\[
    \widetilde f\rtimes_W\widetilde Q,
\]
and applying
\(\Sigma_{\pi_{XZ}}^{\mathsf K}\) under the canonical Frobenius
identifications.

Dually, for an ordinary span
\[
    R:U\rightsquigarrow X,
\]
define the pure right whiskering
\[
    \widehat f\blacktriangleleft R
\]
by pulling to \(U\times X\times Y\), forming the corresponding left
pure-parameter extension, and applying the relevant dependent-sum transport.

\begin{proposition}[Functoriality of pure span whiskering]
\label{prop:functorial-pure-span-whiskering}
For every ordinary span \(Q:Y\rightsquigarrow Z\),
\[
    Q\blacktriangleright(-)
\]
is a functor between the corresponding Kleisli hom-categories of effectful
span cells.  For every ordinary span \(R:U\rightsquigarrow X\),
\[
    (-)\blacktriangleleft R
\]
is likewise a functor.

These functors preserve central cells.
\end{proposition}

\begin{proof}
Each whiskering is a composite of:
\begin{enumerate}
    \item pullback reindexing;
    \item one of the pure-parameter Kleisli functors of
    Proposition~\ref{prop:pure-parameter-kleisli-functors};
    \item a dependent-sum Kleisli functor;
    \item canonical Frobenius and pullback identifications.
\end{enumerate}
Every constituent is functorial.  Hence the composite preserves Kleisli
identities and composition.

Pullback reindexing preserves centrality.  Tensoring with a pure arrow
preserves centrality because pure arrows are central and central arrows are
closed under their common tensor.  Dependent-sum transport preserves
centrality by
Proposition~\ref{prop:fibred-central-arrow-closure}.
\end{proof}

\begin{proposition}[Coherence of pure span whiskering]
\label{prop:coherence-pure-span-whiskering}
There are canonical coherent isomorphisms
\[
    R\blacktriangleright
    (Q\blacktriangleright\widehat f)
    \cong
    (R\circ Q)\blacktriangleright\widehat f,
\]
\[
    (\widehat f\blacktriangleleft Q)
    \blacktriangleleft R
    \cong
    \widehat f\blacktriangleleft(Q\circ R),
\]
and
\[
    (Q\blacktriangleright\widehat f)
    \blacktriangleleft R
    \cong
    Q\blacktriangleright
    (\widehat f\blacktriangleleft R).
\]
They satisfy the pentagon, triangle, and mixed-associativity coherence
equations.
\end{proposition}

\begin{proof}
Pull every cell to the common complete interface product.  Both sides of the
first comparison apply the same effectful cell \(\widehat f\), retain the two
pure span parameters, and then forget the same intermediate interfaces.
The two constructions differ only by:
\begin{enumerate}
    \item reassociation of the slice products;
    \item the order in which the dependent sums are performed;
    \item the canonical pullback associator.
\end{enumerate}
Proposition~\ref{prop:pure-parameter-coherence},
Proposition~\ref{prop:dependent-sum-effect-functoriality},
effectful Beck--Chevalley,
Lemma~\ref{lem:dependent-sum-strength-compatibility}, and Frobenius identify
the two maps.

The right-whiskering comparison is dual.  For mixed associativity, one pure
parameter lies on each side of the single effectful cell.  No exchange of two
computations occurs, and the two composites are identified by mixed
pure-parameter coherence and the canonical dependent-sum comparison.

The coherence diagrams reduce on the common interface product to coherence
for cartesian products, the chosen pullback cleavage, Beck--Chevalley,
Frobenius, and functoriality of dependent-sum transport.
\end{proof}

\begin{corollary}[Pure whiskering]
\label{prop:pure-whiskering-effectful-spans}
If either \(\widehat f\) or \(\widehat g\) is pure, then
\[
    \widehat g\ast_\ell\widehat f
    =
    \widehat g\ast_r\widehat f.
\]
Consequently ordinary span cells act unambiguously on effectful span cells by
left and right whiskering.
\end{corollary}

\begin{proof}
Suppose first that \(\widehat f\) is pure.  Then \(\widetilde f\) is central,
so the two Fubini maps agree after precomposition with
\[
    \widetilde f\times_W\widetilde g.
\]
The dependent-sum transport is the same in both definitions.

If \(\widehat g\) is pure, then \(\widetilde g\) is central.  The symmetry
identity
\[
    d^{\ell}_{P,Q}
    =
    \mathsf K(\tau_{Q,P})
    \circ
    d^{r}_{Q,P}
    \circ
    \tau_{\mathsf KP,\mathsf KQ}
\]
transports centrality from the first coordinate to the second coordinate.
Hence the two Fubini composites again agree before dependent-sum transport.
\end{proof}

\subsection{The interchange obstruction}
\label{subsec:effectful-interchange-obstruction}

For \(P\to Z\), define the diagonal endospan
\[
    D_Z(P)
    :=
    \left(
        Z\xleftarrow{p}P\xrightarrow{p}Z
    \right).
\]
Its object over \(Z\times Z\) is \(\Sigma_{\Delta_Z}P\), where
\[
    \Delta_Z:Z\to Z\times Z
\]
is the diagonal.

For a Kleisli arrow
\[
    f:P\to\mathsf K_ZQ,
\]
define
\[
    D_Z(f)
    :=
    \Sigma_{\Delta_Z}^{\mathsf K}f:
    D_Z(P)\Rrightarrow_{\mathsf K}D_Z(Q).
\]

\begin{lemma}[Diagonal recovery of the Fubini maps]
\label{lem:diagonal-recovery-fubini}
For Kleisli arrows
\[
    f:P\to\mathsf K_ZQ,
    \qquad
    g:R\to\mathsf K_ZS,
\]
there are canonical identifications
\[
    D_Z(R)\circ D_Z(P)
    \cong
    D_Z(P\times_ZR)
\]
and
\[
    D_Z(S)\circ D_Z(Q)
    \cong
    D_Z(Q\times_ZS)
\]
under which
\[
    D_Z(g)\ast_\varepsilon D_Z(f)
    \cong
    D_Z
    \left(
        d^{\varepsilon,Z}_{Q,S}
        \circ
        (f\times_Zg)
    \right),
    \qquad
    \varepsilon\in\{\ell,r\}.
\]
Moreover, \(D_Z\) is faithful on Kleisli arrows.
\end{lemma}

\begin{proof}
The pullback apex of
\[
    D_Z(R)\circ D_Z(P)
\]
is \(P\times_ZR\), giving the first canonical identification, and similarly
for the codomain.

The definition of ordered horizontal composition pulls the two diagonal
cells to the common interface \(Z\times Z\times Z\).  Beck--Chevalley for the
two diagonal pullback squares identifies the resulting reindexed objects
with the pullbacks of \(P,Q,R,S\) over their common \(Z\)-coordinate.
The projection formula of
Lemma~\ref{lem:effectful-projection-formula} then identifies the
dependent-sum transport of the Fubini composite with
\[
    \Sigma_{\Delta_Z}^{\mathsf K}
    \left(
        d^{\varepsilon,Z}_{Q,S}
        \circ
        (f\times_Zg)
    \right).
\]
This is the required diagonal image.

Because \(\Delta_Z\) is monic,
\[
    \Delta_Z^*\Sigma_{\Delta_Z}\cong1.
\]
Effectful Beck--Chevalley therefore gives
\[
    \Delta_Z^{*\mathsf K}
    \Sigma_{\Delta_Z}^{\mathsf K}
    \cong1
\]
on Kleisli arrows, so \(D_Z\) is faithful.
\end{proof}

For Kleisli arrows
\[
    f:P\to\mathsf K_ZQ,
    \qquad
    g:R\to\mathsf K_ZS,
\]
write
\[
    f\otimes_Z^rg
    :=
    d^{r,Z}_{Q,S}\circ(f\times_Zg).
\]
Let
\[
    j_P:=\eta^Z_P
\]
denote the Kleisli identity on \(P\).

\begin{lemma}[Sequential decompositions of the Fubini maps]
\label{lem:sequential-fubini-decompositions}
One has
\[
    (j_Q\otimes_Z^rg)
    \bullet
    (f\otimes_Z^rj_R)
    =
    d^{r,Z}_{Q,S}\circ(f\times_Zg)
\]
and
\[
    (f\otimes_Z^rj_S)
    \bullet
    (j_P\otimes_Z^rg)
    =
    d^{\ell,Z}_{Q,S}\circ(f\times_Zg).
\]
\end{lemma}

\begin{proof}
The first composite expands as
\[
\begin{aligned}
P\times_ZR
&\xrightarrow{f\times1}
\mathsf K_ZQ\times_ZR
\xrightarrow{\operatorname{st}^{r,Z}}
\mathsf K_Z(Q\times_ZR)
\\
&\xrightarrow{\mathsf K_Z(1\times g)}
\mathsf K_Z(Q\times_Z\mathsf K_ZS)
\xrightarrow{\mathsf K_Z(\operatorname{st}^{\ell,Z})}
\mathsf K_Z\mathsf K_Z(Q\times_ZS)
\\
&\xrightarrow{\mu^Z}
\mathsf K_Z(Q\times_ZS).
\end{aligned}
\]
By naturality of the strengths, this is
\[
    d^{r,Z}_{Q,S}\circ(f\times_Zg).
\]

The second composite expands as
\[
\begin{aligned}
P\times_ZR
&\xrightarrow{1\times g}
P\times_Z\mathsf K_ZS
\xrightarrow{\operatorname{st}^{\ell,Z}}
\mathsf K_Z(P\times_ZS)
\\
&\xrightarrow{\mathsf K_Z(f\times1)}
\mathsf K_Z(\mathsf K_ZQ\times_ZS)
\xrightarrow{\mathsf K_Z(\operatorname{st}^{r,Z})}
\mathsf K_Z\mathsf K_Z(Q\times_ZS)
\\
&\xrightarrow{\mu^Z}
\mathsf K_Z(Q\times_ZS),
\end{aligned}
\]
which is
\[
    d^{\ell,Z}_{Q,S}\circ(f\times_Zg).
\]
\end{proof}

\begin{theorem}[Effectful span interchange]
\label{thm:effectful-span-interchange}
The following conditions are equivalent.
\begin{enumerate}
    \item Every slice monad \(\mathsf K_Z\) is commutative:
    \[
        d^{\ell,Z}_{P,Q}
        =
        d^{r,Z}_{P,Q}
    \]
    for all \(Z,P,Q\).

    \item The two ordered horizontal composites agree for all effectful span
    cells, and unrestricted effectful span cells form a bicategory extending
    the ordinary bicategory of spans through the monad unit.

    \item Either ordered horizontal composition is bifunctorial on all
    effectful span cells.
\end{enumerate}
\end{theorem}

\begin{proof}
Suppose first that every \(\mathsf K_Z\) is commutative.  Then the two ordered
horizontal composites agree.  Middle-four interchange follows from
bifunctoriality of the commutative Kleisli tensor, together with
functoriality of reindexing and dependent-sum transport.  Associativity
follows by pulling three cells to the common interface product, applying
associativity of the commutative Fubini map, and then using functoriality of
\(\kappa\).  The unit laws follow from the strength-unit equations.  Pure
cells reproduce ordinary span cells, so the unrestricted cells form the
claimed bicategory.

The implication from the second condition to the third is immediate.

Assume the third condition, and suppose for definiteness that
\(\ast_r\) is bifunctorial.  Fix \(Z\).  By
Lemma~\ref{lem:diagonal-recovery-fubini}, horizontal composition of diagonal
cells recovers the operation \(\otimes_Z^r\).  Since \(D_Z\) preserves
vertical Kleisli composition and is faithful, bifunctoriality of
\(\ast_r\) implies bifunctoriality of \(\otimes_Z^r\).

For arbitrary Kleisli arrows
\[
    f:P\to\mathsf K_ZQ,
    \qquad
    g:R\to\mathsf K_ZS,
\]
bifunctoriality gives
\[
\begin{aligned}
    (j_Q\otimes_Z^rg)
    \bullet
    (f\otimes_Z^rj_R)
    &=
    (j_Q\bullet f)
    \otimes_Z^r
    (g\bullet j_R)
    \\
    &=
    f\otimes_Z^rg,
\end{aligned}
\]
and also
\[
\begin{aligned}
    (f\otimes_Z^rj_S)
    \bullet
    (j_P\otimes_Z^rg)
    &=
    (f\bullet j_P)
    \otimes_Z^r
    (j_S\bullet g)
    \\
    &=
    f\otimes_Z^rg.
\end{aligned}
\]
Lemma~\ref{lem:sequential-fubini-decompositions} therefore yields
\[
    d^{r,Z}_{Q,S}\circ(f\times_Zg)
    =
    d^{\ell,Z}_{Q,S}\circ(f\times_Zg)
\]
for all \(f,g\).

Now take the Kleisli arrows
\[
    \varepsilon_P:
    \mathsf K_ZP\rightsquigarrow P,
    \qquad
    \varepsilon_Q:
    \mathsf K_ZQ\rightsquigarrow Q,
\]
represented by the identity maps
\[
    1_{\mathsf K_ZP},
    \qquad
    1_{\mathsf K_ZQ}.
\]
Then
\[
    \varepsilon_P\times_Z\varepsilon_Q
    =
    1_{\mathsf K_ZP\times_Z\mathsf K_ZQ},
\]
so
\[
    d^{r,Z}_{P,Q}
    =
    d^{\ell,Z}_{P,Q}.
\]
Thus \(\mathsf K_Z\) is commutative.  Since \(Z\) was arbitrary, every slice
monad is commutative.

The argument for \(\ast_\ell\) is symmetric.
\end{proof}

\subsection{The central effectful span bicategory}
\label{subsec:central-effectful-span-bicategory}

\begin{proposition}[Central horizontal composition]
\label{prop:central-horizontal-composition-spans}
If \(\widehat f\) and \(\widehat g\) are central, then
\[
    \widehat g\ast_\ell\widehat f
    =
    \widehat g\ast_r\widehat f.
\]
Their common value is central.
\end{proposition}

\begin{proof}
Centrality is preserved by reindexing, so
\(\widetilde f\) and \(\widetilde g\) are central in the slice over \(W\).
Hence their two Fubini composites agree.

Their common tensor is central by
Proposition~\ref{prop:fibred-central-arrow-closure}, item \(3\).
Dependent-sum transport preserves centrality by item \(5\) of the same
proposition.  Therefore the common horizontal composite is central.
\end{proof}

For central cells, write
\[
    \widehat g\ast\widehat f
\]
for the common horizontal composite.

\begin{theorem}[The central effectful span bicategory]
\label{thm:central-effectful-span-bicategory}
There is a bicategory
\[
    \mathbf{Span}^{\mathrm{cen}}_{\mathsf K}(\mathcal E)
\]
whose:
\begin{enumerate}
    \item objects are objects of \(\mathcal E\);

    \item 1-cells are spans in \(\mathcal E\);

    \item 2-cells are central \(\mathsf K\)-effectful span cells;

    \item vertical composition is fibrewise Kleisli composition;

    \item horizontal composition is the common central composite \(\ast\).
\end{enumerate}
The associator and unitors are the pure effectful images of the canonical
pullback associator and unitors of \(\mathbf{Span}(\mathcal E)\).
\end{theorem}

\begin{proof}
The hom-categories exist by closure of central arrows under Kleisli
composition.  Horizontal composition is well-defined and closed by
Proposition~\ref{prop:central-horizontal-composition-spans}.

For associativity, pull three horizontally composable cells to the common
interface product
\[
    X\times Y\times Z\times T.
\]
Because all three reindexed arrows are central, all bracketings and all
execution orders of their threefold Fubini composite agree.  Both
parenthesizations of horizontal composition are therefore the dependent-sum
transport of the same threefold central tensor.  The remaining comparison is
the canonical pullback associator.

Middle-four interchange follows from bifunctoriality of the monoidal product
in each local centre, followed by functoriality of reindexing and
\(\Sigma^{\mathsf K}\).

The unit laws follow from the strength-unit equations and the pullback
unitors.  Since all coherence 2-cells are pure, the bicategorical coherence
diagrams reduce to the corresponding diagrams in
\(\mathbf{Span}(\mathcal E)\).
\end{proof}

\begin{corollary}[Pure span image]
\label{cor:pure-span-embedding}
The assignment
\[
    f\longmapsto\eta\circ f
\]
extends to an identity-on-objects normal pseudofunctor
\[
    J_{\mathsf K}:
    \mathbf{Span}(\mathcal E)
    \longrightarrow
    \mathbf{Span}^{\mathrm{cen}}_{\mathsf K}(\mathcal E).
\]
\end{corollary}

\section{Effectful Mealy morphisms}
\label{sec:effectful-mealy-morphisms}

Fix internal categories
\[
    \mathbb A=(A_0,A_1,s_A,t_A,e_A,m_A),
    \qquad
    \mathbb B=(B_0,B_1,s_B,t_B,e_B,m_B).
\]
Write
\[
    \mathcal C_{A,B}
    :=
    \mathbf{Span}(\mathcal E)(A_0,B_0)
    \cong
    \mathcal E/(A_0\times B_0).
\]
As in Chapter~2, put
\[
    R_A=-\circ A,
    \qquad
    T_B=B\circ-.
\]
The internal-category structures make \(R_A\) and \(T_B\) monads on
\(\mathcal C_{A,B}\), and Chapter~2 constructs a canonical distributive law
\[
    \lambda:
    R_AT_B
    \Longrightarrow
    T_BR_A.
\]

\subsection{Canonical interaction with input and output}
\label{subsec:canonical-effect-input-output}

The endofunctors \(R_A\) and \(T_B\) are built from pullback and dependent
sum.  The fibred monad therefore supplies canonical transformations moving
an effect through either functor.

For \(P\in\mathcal C_{A,B}\), regard the ordinary identity
\[
    1_{\mathsf K P}:
    \mathsf K P\longrightarrow\mathsf K P
\]
as the Kleisli arrow
\[
    \varepsilon_P:
    \mathsf K P\rightsquigarrow P.
\]

\begin{proposition}[Canonical output and input transformations]
\label{prop:canonical-effect-transformations}
There are canonical natural transformations
\[
    \sigma_B:
    T_B\mathsf K
    \Longrightarrow
    \mathsf KT_B
\]
and
\[
    \rho_A:
    R_A\mathsf K
    \Longrightarrow
    \mathsf KR_A.
\]
The component \(\sigma_{B,P}\) is pure left whiskering of
\(\varepsilon_P\) by the identity cell on \(B\).  The component
\(\rho_{A,P}\) is pure right whiskering of \(\varepsilon_P\) by the identity
cell on \(A\).
\end{proposition}

\begin{proof}
By Proposition~\ref{prop:functorial-pure-span-whiskering}, pure left
whiskering by \(B\) defines a functor on the relevant Kleisli hom-category.
Apply it to
\[
    \varepsilon_P:\mathsf KP\rightsquigarrow P.
\]
The result is an effectful cell
\[
    B\circ\mathsf KP
    \Rrightarrow_{\mathsf K}
    B\circ P,
\]
represented by an ordinary map
\[
    T_B\mathsf KP
    \longrightarrow
    \mathsf KT_BP.
\]
This is \(\sigma_{B,P}\).

Similarly, pure right whiskering by \(A\) sends \(\varepsilon_P\) to an
effectful cell
\[
    \mathsf KP\circ A
    \Rrightarrow_{\mathsf K}
    P\circ A,
\]
represented by
\[
    R_A\mathsf KP
    \longrightarrow
    \mathsf KR_AP.
\]
This is \(\rho_{A,P}\).

Naturality follows from functoriality of the two pure-whiskering functors.
\end{proof}

In internal generalized-element notation, \(\sigma_B\) moves a pure output
arrow into the computation, while \(\rho_A\) moves a pure input arrow into
the computation:
\[
    (\kappa,\beta)
    \longmapsto
    \mathsf K(x\mapsto(x,\beta))(\kappa),
\]
\[
    (a,\kappa)
    \longmapsto
    \mathsf K(x\mapsto(a,x))(\kappa).
\]

\begin{theorem}[Canonical iterated distributive laws]
\label{thm:canonical-iterated-distributive-laws}
The transformations
\[
    \sigma_B:T_B\mathsf K\Rightarrow\mathsf KT_B,
    \qquad
    \rho_A:R_A\mathsf K\Rightarrow\mathsf KR_A
\]
are distributive laws of monads.  Together with
\[
    \lambda:R_AT_B\Rightarrow T_BR_A,
\]
they satisfy, under the standing chosen associativity and pullback
identifications, the Yang--Baxter equation
\[
\begin{aligned}
    (\mathsf K\lambda)
    \circ
    (\rho_AT_B)
    \circ
    (R_A\sigma_B)
    &=
    (\sigma_BR_A)
    \circ
    (T_B\rho_A)
    \circ
    (\lambda\mathsf K).
\end{aligned}
\]
\end{theorem}

\begin{proof}
Pure left whiskering by \(B\) defines a functor
\[
    \widetilde T_B:
    \operatorname{Kl}(\mathsf K)
    \longrightarrow
    \operatorname{Kl}(\mathsf K).
\]
For a Kleisli arrow
\[
    f:P\to\mathsf KQ,
\]
factor \(f\) in the Kleisli category as
\[
    P
    \xrightarrow{\eta_{\mathsf KQ}\circ f}
    \mathsf KQ
    \xrightarrow{\varepsilon_Q}
    Q.
\]
The first factor is pure, and the second is \(\varepsilon_Q\).
Applying pure left whiskering gives the underlying map
\[
    T_BP
    \xrightarrow{T_Bf}
    T_B\mathsf KQ
    \xrightarrow{\sigma_{B,Q}}
    \mathsf KT_BQ.
\]
Thus
\[
    \widetilde T_B(f)
    =
    \sigma_{B,Q}\circ T_Bf.
\]

The unit and multiplication cells of \(\mathbb B\) are pure.  By
Proposition~\ref{prop:coherence-pure-span-whiskering}, their left
whiskerings define a monad structure on \(\widetilde T_B\), and the monad
equations reduce to those of \(\mathbb B\).  Hence \(\widetilde T_B\) is a
lifting of the monad \(T_B\) to \(\operatorname{Kl}(\mathsf K)\).
By the Kleisli lifting--distributive-law correspondence, \(\sigma_B\) is a
distributive law.

The same argument with pure right whiskering by \(A\) gives
\[
    \widetilde R_A(f)
    =
    \rho_{A,Q}\circ R_Af
\]
and proves that \(\rho_A\) is a distributive law.

For the Yang--Baxter equation, fix \(P\).  Both sides are maps
\[
    R_AT_B\mathsf KP
    \longrightarrow
    \mathsf KT_BR_AP.
\]
The left-hand side first left-whiskers
\(\varepsilon_P\) by \(B\), then right-whiskers the result by \(A\), and
finally applies the canonical span associator \(\lambda\).

The right-hand side first applies \(\lambda\), then right-whiskers
\(\varepsilon_P\) by \(A\), and finally left-whiskers by \(B\).

By the mixed-associativity comparison of
Proposition~\ref{prop:coherence-pure-span-whiskering}, both composites
represent the same effectful span cell
\[
    (B\circ\mathsf KP)\circ A
    \Rrightarrow_{\mathsf K}
    B\circ(P\circ A),
\]
obtained by whiskering the single effectful cell
\(\varepsilon_P\) on the left by \(B\) and on the right by \(A\).
Under the standing chosen coherence identifications, this gives the displayed
Yang--Baxter equation.
\end{proof}

\subsection{The effectful output-comparison monad}
\label{subsec:effectful-output-comparison-monad}

\begin{proposition}[Effectful output comparison]
\label{prop:effectful-output-comparison-monad}
The composite
\[
    H_B^{\mathsf K}
    :=
    \mathsf K_{A_0\times B_0}T_B
\]
is a monad on \(\mathcal C_{A,B}\).  Its unit is
\[
    P
    \xrightarrow{\eta^T_P}
    T_BP
    \xrightarrow{\eta^{\mathsf K}_{T_BP}}
    \mathsf KT_BP,
\]
and its multiplication is
\[
\begin{tikzcd}[column sep=large]
    \mathsf KT_B\mathsf KT_BP
        \arrow[r,"{\mathsf K\sigma_BT_BP}"]
    &
    \mathsf K\mathsf KT_BT_BP
        \arrow[r,"{\mu^{\mathsf K}T_BT_B}"]
    &
    \mathsf KT_BT_BP
        \arrow[r,"{\mathsf K\mu^T_P}"]
    &
    \mathsf KT_BP.
\end{tikzcd}
\]
\end{proposition}

\begin{proof}
This is the composite-monad construction for the distributive law
\[
    \sigma_B:T_B\mathsf K\Rightarrow\mathsf KT_B.
\]
\end{proof}

A pure outcome of \(H_B^{\mathsf K}(P)\) is a pair
\[
    (x,\beta)
\]
with
\[
    \beta:q(x)\to b.
\]
The state and output-comparison arrow therefore occur inside the same
computation.

\begin{proposition}[Effectful input distributive law]
\label{prop:effectful-input-distributive-law}
Define
\[
    \chi:
    R_AH_B^{\mathsf K}
    \Longrightarrow
    H_B^{\mathsf K}R_A
\]
by
\[
    \chi
    :=
    (\mathsf K\lambda)
    \circ
    (\rho_AT_B).
\]
Then \(\chi\) is a distributive law of the monad \(R_A\) over the monad
\(H_B^{\mathsf K}\).  Hence \(R_A\) lifts to a monad
\[
    \overline R_A^{\mathsf K}
\]
on
\[
    \operatorname{Kl}(H_B^{\mathsf K}).
\]
\end{proposition}

\begin{proof}
Apply the iterated distributive-law theorem to the three monads
\[
    R_A,
    \qquad
    T_B,
    \qquad
    \mathsf K
\]
and the Yang--Baxter equation of
Theorem~\ref{thm:canonical-iterated-distributive-laws}.
\end{proof}

\subsection{Definition and algebra presentation}
\label{subsec:effectful-mealy-definition}

\begin{definition}[Effectful Mealy morphism]
\label{def:effectful-mealy-morphism}
A \textbf{\(\mathsf K\)-effectful Mealy morphism}
\[
    (P,\widehat\Phi):
    \mathbb A\rightsquigarrow_{\mathsf K}\mathbb B
\]
consists of a carrier span
\[
    A_0\xleftarrow{p}P\xrightarrow{q}B_0
\]
and an algebra structure for
\[
    \overline R_A^{\mathsf K}
\]
on \(P\) in
\[
    \operatorname{Kl}(H_B^{\mathsf K}).
\]

Equivalently, it consists of a map
\[
    \widehat\Phi:
    P\circ A
    \longrightarrow
    \mathsf K_{A_0\times B_0}(B\circ P)
\]
satisfying
\begin{align}
    \widehat\Phi\circ\eta^R_P
    &=
    \eta^H_P,
    \label{eq:effectful-mealy-unit}
    \\
    \widehat\Phi\circ\mu^R_P
    &=
    \mu^H_P
    \circ
    H_B^{\mathsf K}(\widehat\Phi)
    \circ
    \chi_P
    \circ
    R_A(\widehat\Phi).
    \label{eq:effectful-mealy-multiplication}
\end{align}
\end{definition}

\begin{definition}[Central effectful Mealy morphism]
\label{def:central-effectful-mealy-morphism}
An effectful Mealy morphism is \textbf{central} when its structure map
\[
    \widehat\Phi:
    P\circ A
    \longrightarrow
    \mathsf K(B\circ P)
\]
is a central Kleisli arrow in
\(\mathcal E/(A_0\times B_0)\).
\end{definition}

\begin{theorem}[Effectful Mealy algebra theorem]
\label{thm:effectful-mealy-algebra-theorem}
For a fixed carrier span \(P:A_0\rightsquigarrow B_0\), giving an effectful
Mealy structure on \(P\) is equivalently giving an algebra for
\[
    \overline R_A^{\mathsf K}
\]
on \(P\) in
\[
    \operatorname{Kl}(H_B^{\mathsf K}).
\]

If the structure map is central, this algebra is equivalently an oriented
monad morphism from \(\mathbb A\) to \(\mathbb B\) in the bicategory
\[
    \mathbf{Span}^{\mathrm{cen}}_{\mathsf K}(\mathcal E),
\]
with carrier \(P\) and structure cell
\[
    \widehat\Phi:
    P\circ A
    \Rrightarrow_{\mathsf K}
    B\circ P.
\]
\end{theorem}

\begin{proof}
The first equivalence is the definition of an algebra for the lifted monad:
an algebra arrow in \(\operatorname{Kl}(H_B^{\mathsf K})\) is an ordinary map
\[
    R_AP
    \longrightarrow
    H_B^{\mathsf K}P
\]
satisfying the two displayed algebra equations.

If \(\widehat\Phi\) is central, it is a 2-cell in
\(\mathbf{Span}^{\mathrm{cen}}_{\mathsf K}(\mathcal E)\).  The oriented monad
unit diagram is precisely
Equation~\eqref{eq:effectful-mealy-unit}.  Expanding the oriented monad
multiplication diagram using the distributive law \(\chi\) gives precisely
Equation~\eqref{eq:effectful-mealy-multiplication}.  Conversely, those two
equations supply the oriented monad compatibility diagrams.
\end{proof}

\subsection{Internal component laws}
\label{subsec:internal-effectful-mealy-laws}

For generalized elements \(u:A_0\) and \(b:B_0\), write
\[
    \operatorname{Out}_P(u,b)
\]
for the fibre of
\[
    B\circ P\longrightarrow A_0\times B_0
\]
over \((u,b)\).  Internally its elements are pairs
\[
    (x',\beta)
\]
satisfying
\[
    p(x')=u,
    \qquad
    s_B(\beta)=q(x'),
    \qquad
    t_B(\beta)=b.
\]

For an admissible pair
\[
    a:u\to v,
    \qquad
    p(x)=v,
\]
the transition has type
\[
    \widehat\Phi(a,x)
    \in
    \mathsf K
    \bigl(
        \operatorname{Out}_P(u,q(x))
    \bigr).
\]

Write
\[
    \operatorname{bind}_{\mathsf K}(\kappa,f)
    :=
    \mu\circ\mathsf K(f)(\kappa).
\]

\begin{proposition}[Effectful Mealy component laws]
\label{prop:effectful-mealy-component-laws}
Equations~\eqref{eq:effectful-mealy-unit} and
\eqref{eq:effectful-mealy-multiplication} are equivalent to the following
internal equations.

The unit law is
\[
    \widehat\Phi(e_A(px),x)
    =
    \eta^{\mathsf K}
    \bigl(
        x,e_B(qx)
    \bigr).
\]

For internally composable arrows
\[
    a:u\to v,
    \qquad
    b:v\to w,
\]
and \(p(x)=w\), the multiplication law is
\[
\begin{aligned}
    \widehat\Phi(b\circ a,x)
    ={}&
    \operatorname{bind}_{\mathsf K}
    \Bigl(
        \widehat\Phi(b,x),
        (x_b,\beta_b)
        \longmapsto
    \\
    &\qquad
        \operatorname{bind}_{\mathsf K}
        \Bigl(
            \widehat\Phi(a,x_b),
            (x_a,\beta_a)
            \longmapsto
            \eta^{\mathsf K}
            \bigl(
                x_a,
                \beta_b\circ\beta_a
            \bigr)
        \Bigr)
    \Bigr).
\end{aligned}
\]
\end{proposition}

\begin{proof}
The unit of \(R_A\) inserts the identity input arrow.  The unit of
\(H_B^{\mathsf K}\) inserts the identity output arrow and then applies
\(\eta^{\mathsf K}\), giving the first equation.

For multiplication, target-to-source execution first processes
\[
    b:v\to w
\]
at \(x\), producing
\[
    (x_b,\beta_b),
    \qquad
    \beta_b:q(x_b)\to q(x).
\]
It then processes
\[
    a:u\to v
\]
at \(x_b\), producing
\[
    (x_a,\beta_a),
    \qquad
    \beta_a:q(x_a)\to q(x_b).
\]
The output-comparison multiplication composes these arrows to
\[
    \beta_b\circ\beta_a:
    q(x_a)\to q(x),
\]
and the multiplication of \(\mathsf K\) combines the nested computations.
\end{proof}

\subsection{Pure effectful Mealy morphisms}
\label{subsec:pure-effectful-mealy-morphisms}

\begin{proposition}[Pure effectful image]
\label{prop:pure-effectful-image}
Every deterministic Mealy morphism
\[
    (P,\Phi):
    \mathbb A\rightsquigarrow\mathbb B
\]
determines a central effectful Mealy morphism with the same carrier by
\[
    \widehat\Phi
    :=
    \eta^{\mathsf K}_{B\circ P}\circ\Phi.
\]
\end{proposition}

\begin{proof}
Apply the pure pseudofunctor
\[
    J_{\mathsf K}:
    \mathbf{Span}(\mathcal E)
    \to
    \mathbf{Span}^{\mathrm{cen}}_{\mathsf K}(\mathcal E)
\]
to the oriented monad morphism \((P,\Phi)\).  Pure Kleisli arrows are
central.
\end{proof}

\section{Effectful Mealy simulations}
\label{sec:effectful-mealy-simulations}

Before defining simulations, we record closure under the Kleisli composition
of the output-comparison monad.

\begin{proposition}[Central output-comparison composition]
\label{prop:central-output-comparison-composition}
Let
\[
    f:P\to\mathsf KT_BQ,
    \qquad
    g:Q\to\mathsf KT_BR
\]
be central as \(\mathsf K\)-Kleisli arrows.  Their
\(H_B^{\mathsf K}\)-Kleisli composite
\[
    g\mathbin{\bullet_H}f
    :=
    \mu^H_R
    \circ
    H_B^{\mathsf K}(g)
    \circ
    f
\]
is central.  The \(H_B^{\mathsf K}\)-Kleisli identities are central.
\end{proposition}

\begin{proof}
Expand the composite:
\[
\begin{aligned}
P
&\xrightarrow{f}
\mathsf KT_BQ
\xrightarrow{\mathsf K(T_Bg)}
\mathsf KT_B\mathsf KT_BR
\\
&\xrightarrow{\mathsf K(\sigma_BT_BR)}
\mathsf K\mathsf KT_BT_BR
\xrightarrow{\mu^{\mathsf K}}
\mathsf KT_BT_BR
\xrightarrow{\mathsf K(\mu^T_R)}
\mathsf KT_BR.
\end{aligned}
\]
Pure left whiskering by \(B\) preserves centrality by
Proposition~\ref{prop:functorial-pure-span-whiskering}.  Thus the part
induced by \(g\) is central.  The displayed map is then a Kleisli composite
of central arrows followed by a pure map, so it is central.  The identity is
the composite of pure units and is therefore central.
\end{proof}

Let
\[
    (P,\widehat\Phi),
    (Q,\widehat\Psi):
    \mathbb A\rightsquigarrow_{\mathsf K}\mathbb B
\]
be effectful Mealy morphisms.

\begin{definition}[Effectful Mealy simulation]
\label{def:effectful-mealy-simulation}
An \textbf{effectful Mealy simulation}
\[
    (P,\widehat\Phi)
    \Rightarrow
    (Q,\widehat\Psi)
\]
is a morphism
\[
    (Q,\widehat\Psi)
    \longrightarrow
    (P,\widehat\Phi)
\]
of \(\overline R_A^{\mathsf K}\)-algebras in
\(\operatorname{Kl}(H_B^{\mathsf K})\).

Equivalently, it is a map
\[
    \widehat\Theta:
    Q\longrightarrow\mathsf K(B\circ P)
\]
satisfying
\[
\begin{aligned}
&
\mu^H_P
\circ
H_B^{\mathsf K}(\widehat\Theta)
\circ
\widehat\Psi
\\
&\qquad=
\mu^H_P
\circ
H_B^{\mathsf K}(\widehat\Phi)
\circ
\chi_P
\circ
R_A(\widehat\Theta).
\end{aligned}
\tag{\(\mathrm{Sim}_{\mathsf K}\)}
\]
The simulation is \textbf{central} when \(\widehat\Theta\) is central as a
\(\mathsf K\)-Kleisli arrow.
\end{definition}

The displayed simulation direction is opposite to the underlying
algebra-morphism direction, exactly as in Chapter~2.

\begin{proposition}[Component simulation equation]
\label{prop:effectful-simulation-component-equation}
Internally, write
\[
    \widehat\Theta(y)
    \in
    \mathsf K
    \left\{
        (x,\alpha)
        \,\middle|\,
        \alpha:q_P(x)\to q_Q(y)
    \right\}.
\]
Then \(\widehat\Theta\) is an effectful Mealy simulation if and only if, for
every admissible \((a,y)\),
\[
\begin{aligned}
&
\operatorname{bind}_{\mathsf K}
\Bigl(
    \widehat\Theta(y),
    (x_0,\alpha_0)
    \longmapsto
    \operatorname{bind}_{\mathsf K}
    \Bigl(
        \widehat\Phi(a,x_0),
        (x,\beta)
        \longmapsto
        \eta^{\mathsf K}
        \bigl(
            x,
            \alpha_0\circ\beta
        \bigr)
    \Bigr)
\Bigr)
\\
&\qquad=
\operatorname{bind}_{\mathsf K}
\Bigl(
    \widehat\Psi(a,y),
    (y',\gamma)
    \longmapsto
    \operatorname{bind}_{\mathsf K}
    \Bigl(
        \widehat\Theta(y'),
        (x,\alpha)
        \longmapsto
        \eta^{\mathsf K}
        \bigl(
            x,
            \gamma\circ\alpha
        \bigr)
    \Bigr)
\Bigr).
\end{aligned}
\]
\end{proposition}

\begin{proof}
The left-hand side first compares \(y\) with an effectful
\(P\)-state--output pair \((x_0,\alpha_0)\), then executes \(P\) along \(a\),
and composes
\[
    q_P(x)
    \xrightarrow{\beta}
    q_P(x_0)
    \xrightarrow{\alpha_0}
    q_Q(y).
\]
The right-hand side first executes \(Q\), then compares the updated state,
and composes
\[
    q_P(x)
    \xrightarrow{\alpha}
    q_Q(y')
    \xrightarrow{\gamma}
    q_Q(y).
\]
Expanding \((\mathrm{Sim}_{\mathsf K})\) gives exactly the displayed equality.
\end{proof}

\begin{definition}[Pure effectful simulation]
\label{def:pure-effectful-simulation}
A simulation is \textbf{pure} when it has the form
\[
    \widehat\Theta
    =
    \eta^{\mathsf K}_{B\circ P}\circ\Theta
\]
for an ordinary output-comparison cell
\[
    \Theta:Q\longrightarrow B\circ P.
\]
\end{definition}

Writing
\[
    \Theta(y)
    =
    (\theta y,\alpha_y),
\]
the component simulation equation reduces to
\[
\begin{aligned}
&
\operatorname{bind}_{\mathsf K}
\Bigl(
    \widehat\Phi(a,\theta y),
    (x,\beta)
    \longmapsto
    \eta^{\mathsf K}
    \bigl(
        x,
        \alpha_y\circ\beta
    \bigr)
\Bigr)
\\
&\qquad=
\operatorname{bind}_{\mathsf K}
\Bigl(
    \widehat\Psi(a,y),
    (y',\gamma)
    \longmapsto
    \eta^{\mathsf K}
    \bigl(
        \theta y',
        \gamma\circ\alpha_{y'}
    \bigr)
\Bigr).
\end{aligned}
\]

\begin{proposition}[Vertical composition of simulations]
\label{prop:vertical-composition-effectful-simulations}
Let
\[
    \widehat\Theta:
    (P,\widehat\Phi)
    \Rightarrow
    (Q,\widehat\Psi)
\]
and
\[
    \widehat\Lambda:
    (Q,\widehat\Psi)
    \Rightarrow
    (R,\widehat\Xi)
\]
be effectful Mealy simulations.  Their displayed vertical composite is
written
\[
    \widehat\Lambda\circ_v\widehat\Theta:
    (P,\widehat\Phi)
    \Rightarrow
    (R,\widehat\Xi).
\]
Its representing \(H_B^{\mathsf K}\)-Kleisli arrow is
\[
    \widehat\Theta\mathbin{\bullet_H}\widehat\Lambda.
\]
Internally,
\[
\begin{aligned}
    (\widehat\Lambda\circ_v\widehat\Theta)(z)
    ={}&
    \operatorname{bind}_{\mathsf K}
    \Bigl(
        \widehat\Lambda(z),
        (y,\beta)
        \longmapsto
    \\
    &\qquad
        \operatorname{bind}_{\mathsf K}
        \Bigl(
            \widehat\Theta(y),
            (x,\alpha)
            \longmapsto
            \eta^{\mathsf K}
            \bigl(
                x,
                \beta\circ\alpha
            \bigr)
        \Bigr)
    \Bigr).
\end{aligned}
\]
Central simulations are closed under this composition.
\end{proposition}

\begin{proof}
Displayed simulations are opposite to the corresponding algebra morphisms.
Thus their displayed vertical composite is represented by first applying
\(\widehat\Lambda\) and then applying \(\widehat\Theta\) in
\(\operatorname{Kl}(H_B^{\mathsf K})\), namely by
\[
    \widehat\Theta\mathbin{\bullet_H}\widehat\Lambda.
\]
The component formula is the multiplication of \(H_B^{\mathsf K}\).
Centrality follows from
Proposition~\ref{prop:central-output-comparison-composition}.
\end{proof}

\section{Pipeline composition and the effectful Mealy bicategory}
\label{sec:effectful-pipeline-bicategory}

Let
\[
    (P,\widehat\Phi_P):
    \mathbb A\rightsquigarrow_{\mathsf K}\mathbb B,
    \qquad
    (Q,\widehat\Phi_Q):
    \mathbb B\rightsquigarrow_{\mathsf K}\mathbb C
\]
be central effectful Mealy morphisms.

\subsection{Pipeline composition}
\label{subsec:central-effectful-pipeline-composition}

The composite carrier is the ordinary pullback span
\[
    A_0
    \xleftarrow{}
    P\times_{B_0}Q
    \xrightarrow{}
    C_0.
\]

\begin{definition}[Effectful pipeline transition]
\label{def:effectful-pipeline-transition}
The transition of the composite
\[
    Q\circ P:
    \mathbb A\rightsquigarrow_{\mathsf K}\mathbb C
\]
is the horizontal composite of the two oriented monad cells in
\(\mathbf{Span}^{\mathrm{cen}}_{\mathsf K}(\mathcal E)\).

Internally,
\[
\begin{aligned}
&
\widehat\Phi_{Q\circ P}(a,(x,y))
\\
&\quad=
\operatorname{bind}_{\mathsf K}
\Bigl(
    \widehat\Phi_P(a,x),
    (x',\beta)
    \longmapsto
    \operatorname{bind}_{\mathsf K}
    \Bigl(
        \widehat\Phi_Q(\beta,y),
        (y',\gamma)
        \longmapsto
        \eta^{\mathsf K}
        \bigl(
            (x',y'),
            \gamma
        \bigr)
    \Bigr)
\Bigr).
\end{aligned}
\]
\end{definition}

If
\[
    \beta:q_P(x')\to q_P(x)=p_Q(y),
\]
then \(\beta\) is admissible at \(y\).  Every outcome
\[
    (y',\gamma)
\]
of the second transition satisfies
\[
    p_Q(y')=q_P(x').
\]
Hence
\[
    (x',y')\in P\times_{B_0}Q.
\]

\begin{theorem}[Effectful pipeline composition]
\label{thm:effectful-pipeline-composition}
Definition~\ref{def:effectful-pipeline-transition} defines a central
effectful Mealy morphism
\[
    Q\circ P:
    \mathbb A\rightsquigarrow_{\mathsf K}\mathbb C.
\]
\end{theorem}

\begin{proof}
By Theorem~\ref{thm:effectful-mealy-algebra-theorem}, the two machines are
oriented monad morphisms in
\(\mathbf{Span}^{\mathrm{cen}}_{\mathsf K}(\mathcal E)\).  Oriented monad
morphisms compose in every bicategory.  Their pasted structure cell has
carrier \(Q\circ P\) and expands to the bind formula of
Definition~\ref{def:effectful-pipeline-transition}.  The oriented monad
unit and multiplication diagrams for the paste follow from those of the two
factors and bicategorical interchange.

The composite structure cell is central by
Proposition~\ref{prop:central-horizontal-composition-spans}.
\end{proof}

\begin{corollary}[Associativity and identities]
\label{cor:effectful-pipeline-associativity}
Effectful pipeline composition is associative up to the canonical pullback
associator.  The pure images of the deterministic identity Mealy morphisms
are left and right units.
\end{corollary}

\subsection{Horizontal composition of simulations}
\label{subsec:horizontal-composition-effectful-simulations}

Let
\[
    \widehat\Theta:
    (P,\widehat\Phi_P)
    \Rightarrow
    (P',\widehat\Phi_{P'})
\]
be a central simulation from \(\mathbb A\) to \(\mathbb B\), and let
\[
    \widehat\Lambda:
    (Q,\widehat\Phi_Q)
    \Rightarrow
    (Q',\widehat\Phi_{Q'})
\]
be a central simulation from \(\mathbb B\) to \(\mathbb C\).

Internally, write
\[
    \widehat\Theta(x')
    \in
    \mathsf K
    \left\{
        (x,\alpha)
        \,\middle|\,
        \alpha:q_P(x)\to q_{P'}(x')
    \right\},
\]
\[
    \widehat\Lambda(y')
    \in
    \mathsf K
    \left\{
        (y,\beta)
        \,\middle|\,
        \beta:q_Q(y)\to q_{Q'}(y')
    \right\}.
\]

\begin{definition}[Horizontal composite of simulations]
\label{def:horizontal-composite-effectful-simulations}
The horizontal composite
\[
    \widehat\Lambda\ast_h\widehat\Theta:
    Q\circ P
    \Rightarrow
    Q'\circ P'
\]
is represented internally by
\[
\begin{aligned}
&
(\widehat\Lambda\ast_h\widehat\Theta)(x',y')
\\
&=
\operatorname{bind}_{\mathsf K}
\Bigl(
    \widehat\Theta(x'),
    (x,\alpha)
    \longmapsto
    \operatorname{bind}_{\mathsf K}
    \Bigl(
        \widehat\Lambda(y'),
        (y,\beta)
        \longmapsto
\\
&\hspace{8em}
        \operatorname{bind}_{\mathsf K}
        \Bigl(
            \widehat\Phi_Q(\alpha,y),
            (y_1,\gamma)
            \longmapsto
            \eta^{\mathsf K}
            \bigl(
                (x,y_1),
                \beta\circ\gamma
            \bigr)
        \Bigr)
    \Bigr)
\Bigr).
\end{aligned}
\]
\end{definition}

The arrow
\[
    \alpha:q_P(x)\to q_{P'}(x')=p_{Q'}(y')=p_Q(y)
\]
is executed by the downstream machine \(Q\).  Its updated state \(y_1\)
satisfies
\[
    p_Q(y_1)=q_P(x),
\]
so
\[
    (x,y_1)\in P\times_{B_0}Q.
\]

\begin{proposition}[Horizontal simulation law]
\label{prop:horizontal-effectful-simulation-law}
The map of
Definition~\ref{def:horizontal-composite-effectful-simulations} is a central
effectful Mealy simulation.  Horizontal composition preserves vertical
identities and satisfies middle-four interchange with respect to
\(\circ_v\).
\end{proposition}

\begin{proof}
Paste the comparison cell \(\widehat\Theta\), the transition cell of \(Q\)
evaluated on its output comparison, and the comparison cell
\(\widehat\Lambda\) in
\(\mathbf{Span}^{\mathrm{cen}}_{\mathsf K}(\mathcal E)\).  The component of
the resulting paste is exactly the displayed bind expression.

The pasted cell satisfies the simulation equation because it is the
horizontal composite of algebra-morphism cells between oriented monad
morphisms.  Equivalently, expanding both sides and applying the simulation
equations for \(\widehat\Theta\) and \(\widehat\Lambda\), followed by the
Mealy multiplication equation for \(Q\), yields the same central
four-computation paste.

The component is built from reindexing, central tensor, Kleisli composition,
pure whiskering, and dependent-sum transport.  Each operation preserves
centrality.

Preservation of identities and middle-four interchange are inherited from
the corresponding laws in
\(\mathbf{Span}^{\mathrm{cen}}_{\mathsf K}(\mathcal E)\), with the displayed
simulation direction reversed relative to algebra morphisms.
\end{proof}

\begin{theorem}[The central effectful Mealy bicategory]
\label{thm:central-effectful-mealy-bicategory}
There is a bicategory
\[
    \mathbf{Mly}^{\mathrm{cen}}_{\mathsf K}(\mathcal E)
\]
whose:
\begin{enumerate}
    \item objects are internal categories in \(\mathcal E\);

    \item 1-cells are central \(\mathsf K\)-effectful Mealy morphisms;

    \item 2-cells are central effectful Mealy simulations.
\end{enumerate}
\end{theorem}

\begin{proof}
For fixed \(\mathbb A,\mathbb B\), the hom-category is the opposite of the
category of \(\overline R_A^{\mathsf K}\)-algebras and central algebra
morphisms in \(\operatorname{Kl}(H_B^{\mathsf K})\).  Vertical closure is
Proposition~\ref{prop:central-output-comparison-composition}.  Horizontal
composition of 1-cells is
Theorem~\ref{thm:effectful-pipeline-composition}, and horizontal composition
of 2-cells is
Proposition~\ref{prop:horizontal-effectful-simulation-law}.

The associator and unitors are induced by the pure pullback associator and
unitors.  Their coherence follows from the corresponding coherence in
\(\mathbf{Span}^{\mathrm{cen}}_{\mathsf K}(\mathcal E)\).
\end{proof}

\begin{theorem}[Deterministic image]
\label{thm:deterministic-effectful-pseudofunctor}
The pure construction extends to a normal pseudofunctor
\[
    J_{\mathsf K}^{\mathrm{Mly}}:
    \mathbf{Mly}(\mathcal E)
    \longrightarrow
    \mathbf{Mly}^{\mathrm{cen}}_{\mathsf K}(\mathcal E).
\]
It is the identity on internal categories and carriers and applies the monad
unit to transition and simulation cells.
\end{theorem}

\begin{proof}
Apply the pure pseudofunctor
\[
    J_{\mathsf K}:
    \mathbf{Span}(\mathcal E)
    \to
    \mathbf{Span}^{\mathrm{cen}}_{\mathsf K}(\mathcal E)
\]
to internal-category monads, oriented monad morphisms, and their simulation
cells.  Preservation of horizontal composition follows from the
strength-unit equations.  Identity cells are sent to Kleisli identities, so
the pseudofunctor is normal.
\end{proof}

No local faithfulness assertion is made: the monad unit need not reflect
equality of ordinary span cells.

\section{Strict state presentations}
\label{sec:effectful-strict-state-presentations}

Simulation 2-cells compare machines while allowing output-comparison arrows.
For presentation changes that act only on the state carrier, the relevant
maps are stricter.

\begin{definition}[Strict state homomorphism]
\label{def:strict-effectful-state-homomorphism}
Let
\[
    (P,\widehat\Phi),
    (Q,\widehat\Psi):
    \mathbb A\rightsquigarrow_{\mathsf K}\mathbb B
\]
be central effectful Mealy morphisms.  A \textbf{strict state homomorphism}
\[
    h:
    (P,\widehat\Phi)
    \longrightarrow
    (Q,\widehat\Psi)
\]
is a map of carrier spans
\[
    h:P\longrightarrow Q
\]
satisfying
\[
    \mathsf K(B\circ h)
    \circ
    \widehat\Phi
    =
    \widehat\Psi
    \circ
    (h\circ A).
\]
\end{definition}

Internally,
\[
    \mathsf K
    \bigl(
        (x',\beta)
        \mapsto
        (hx',\beta)
    \bigr)
    \bigl(
        \widehat\Phi(a,x)
    \bigr)
    =
    \widehat\Psi(a,hx).
\]

A strict state homomorphism \(h:P\to Q\) determines a pure effectful
simulation in the opposite displayed direction
\[
    (Q,\widehat\Psi)
    \Rightarrow
    (P,\widehat\Phi)
\]
whose output-comparison component is the identity.

\begin{proposition}[Composition of strict state homomorphisms]
\label{prop:strict-homomorphism-composition}
Strict state homomorphisms form a category for every fixed pair
\((\mathbb A,\mathbb B)\).  Pipeline composition sends
\[
    h:P\to P',
    \qquad
    k:Q\to Q'
\]
to the pullback map
\[
    h\times k:
    P\times_{B_0}Q
    \longrightarrow
    P'\times_{B_0}Q'.
\]
\end{proposition}

\begin{proof}
Identity maps satisfy the strict transition equation.  If
\[
    P\xrightarrow{h}Q\xrightarrow{k}R
\]
are strict, then
\[
\begin{aligned}
\mathsf K(B\circ kh)\widehat\Phi
&=
\mathsf K(B\circ k)
\mathsf K(B\circ h)
\widehat\Phi
\\
&=
\mathsf K(B\circ k)
\widehat\Psi(h\circ A)
\\
&=
\widehat\Xi(k\circ A)(h\circ A)
\\
&=
\widehat\Xi((kh)\circ A).
\end{aligned}
\]
Thus strict maps compose.

For pipeline composition, apply \(h\) to the updated upstream state and
\(k\) to the updated downstream state.  The two strict transition equations
rewrite the two stages of the bind expression, and naturality of Kleisli
extension identifies the resulting composite with the transition of
\(P'\) followed by that of \(Q'\).
\end{proof}

\begin{theorem}[The bicategory of strict presentations]
\label{thm:strict-presentation-bicategory}
There is a bicategory
\[
    \mathbf{Pres}^{\mathrm{cen}}_{\mathsf K}(\mathcal E)
\]
whose objects and 1-cells are those of
\(\mathbf{Mly}^{\mathrm{cen}}_{\mathsf K}(\mathcal E)\), and whose 2-cells are
strict state homomorphisms in their covariant carrier direction.
\end{theorem}

\begin{proof}
Vertical composition is
Proposition~\ref{prop:strict-homomorphism-composition}.  Horizontal
composition is the pullback map \(h\times k\).  Functoriality of pullback
gives middle-four interchange.

The canonical pullback associator
\[
    (P\times_{B_0}Q)\times_{C_0}R
    \xrightarrow{\cong}
    P\times_{B_0}(Q\times_{C_0}R)
\]
preserves all carrier projections.  Both sides execute the same three
transition cells in the same order, so the associator satisfies the strict
transition equation.  The same argument applies to the pullback unitors.
Their coherence is inherited from the ordinary span bicategory.
\end{proof}

\subsection{Collapse of unrestricted presentation connectedness}
\label{subsec:collapse-unrestricted-state-connectedness}

The connected-component quotient of
\(\mathbf{Pres}^{\mathrm{cen}}_{\mathsf K}(\mathcal E)\) is too coarse under
the conditions needed for feedback.

\begin{proposition}[Empty effectful presentation]
\label{prop:empty-effectful-presentation}
Suppose that \(\mathcal E\) has a pullback-stable initial object \(0\).
For every pair of internal categories \(\mathbb A,\mathbb B\), the empty
carrier span
\[
    A_0\longleftarrow0\longrightarrow B_0
\]
carries a unique central effectful Mealy structure
\[
    \mathbf 0_{\mathbb A,\mathbb B}:
    \mathbb A\rightsquigarrow_{\mathsf K}\mathbb B.
\]
Moreover, there is a unique strict state homomorphism
\[
    \mathbf 0_{\mathbb A,\mathbb B}
    \longrightarrow
    (P,\widehat\Phi)
\]
for every central effectful Mealy morphism
\[
    (P,\widehat\Phi):
    \mathbb A\rightsquigarrow_{\mathsf K}\mathbb B.
\]
\end{proposition}

\begin{proof}
Pullback stability of \(0\) gives
\[
    0\circ A\cong0,
    \qquad
    B\circ0\cong0.
\]
There is therefore a unique map
\[
    \widehat\Phi_0:
    0\longrightarrow
    \mathsf K_{A_0\times B_0}(0).
\]
The Mealy unit and multiplication equations have initial domain, so they
hold.  After arbitrary reindexing, the domain remains initial, and the two
centrality composites are maps out of an initial object.  Hence
\(\widehat\Phi_0\) is central.

For any carrier \(P\), there is a unique span map
\[
    0\to P.
\]
Its strict transition equation is an equality between two maps out of an
initial object, so it is strict.
\end{proof}

\begin{theorem}[Degeneracy of unrestricted strict-map connectedness]
\label{thm:degeneracy-unrestricted-strict-uniformity}
Suppose that \(\mathcal E\) has a pullback-stable initial object.  Then every
two central effectful Mealy morphisms with the same source and target are
connected by a finite zigzag of strict state homomorphisms.

Consequently, the connected-component quotient of
\[
    \mathbf{Pres}^{\mathrm{cen}}_{\mathsf K}
    (\mathbb A,\mathbb B)
\]
is a singleton for every \(\mathbb A,\mathbb B\).
\end{theorem}

\begin{proof}
For any two presentations \(P,Q\),
Proposition~\ref{prop:empty-effectful-presentation} gives a zigzag
\[
    P
    \longleftarrow
    \mathbf0_{\mathbb A,\mathbb B}
    \longrightarrow
    Q.
\]
Thus every pair lies in the same connected component.
\end{proof}

In particular, invariance under all strict state homomorphisms cannot produce
a nontrivial unpointed state-hiding semantics in a feedback span base.

\subsection{Strict state retractions}
\label{subsec:strict-state-retractions}

\begin{definition}[Strict state retraction]
\label{def:strict-state-retraction}
A strict state homomorphism
\[
    h:
    (P,\widehat\Phi)
    \longrightarrow
    (Q,\widehat\Psi)
\]
is a \textbf{strict state retraction} if there exists a strict state
homomorphism
\[
    s:
    (Q,\widehat\Psi)
    \longrightarrow
    (P,\widehat\Phi)
\]
such that
\[
    h\circ s=1_Q.
\]
The map \(s\) is called a \textbf{strict section} of \(h\).
\end{definition}

\begin{proposition}[Closure of strict state retractions]
\label{prop:closure-strict-state-retractions}
Strict state retractions are closed under identity, vertical composition, and
pipeline composition.
\end{proposition}

\begin{proof}
Every identity is a retraction with itself as section.

Let
\[
    P\xrightarrow{h}Q\xrightarrow{k}R
\]
be strict state retractions with strict sections
\[
    s:Q\to P,
    \qquad
    t:R\to Q.
\]
Then \(s\circ t\) is strict and
\[
    (k\circ h)\circ(s\circ t)
    =
    k\circ(h\circ s)\circ t
    =
    k\circ t
    =
    1_R.
\]
Thus \(k\circ h\) is a strict state retraction.

Now let
\[
    h:P\to P',
    \qquad
    k:Q\to Q'
\]
be strict state retractions with strict sections \(s\) and \(t\).
Their horizontal composite is
\[
    h\times k:
    P\times_{B_0}Q
    \longrightarrow
    P'\times_{B_0}Q'.
\]
It has strict section
\[
    s\times t:
    P'\times_{B_0}Q'
    \longrightarrow
    P\times_{B_0}Q,
\]
and
\[
    (h\times k)\circ(s\times t)
    =
    (h\circ s)\times(k\circ t)
    =
    1.
\]
\end{proof}

\begin{theorem}[The bicategory of retractive presentations]
\label{thm:retractive-presentation-bicategory}
There is a locally wide sub-bicategory
\[
    \mathbf{RetPres}^{\mathrm{cen}}_{\mathsf K}(\mathcal E)
    \subseteq
    \mathbf{Pres}^{\mathrm{cen}}_{\mathsf K}(\mathcal E)
\]
with the same objects and 1-cells, and whose 2-cells are the strict state
retractions.
\end{theorem}

\begin{proof}
Closure under vertical and horizontal composition is
Proposition~\ref{prop:closure-strict-state-retractions}.  The associator and
unitors of
\(\mathbf{Pres}^{\mathrm{cen}}_{\mathsf K}(\mathcal E)\)
are strict state isomorphisms, hence strict state retractions.  The
bicategorical coherence equations are inherited from the strict-presentation
bicategory.
\end{proof}

\begin{definition}[Retractive presentation equivalence]
\label{def:retractive-presentation-equivalence}
Two central effectful Mealy morphisms with the same source and target are
\textbf{retractively presentation-equivalent}, written
\[
    (P,\widehat\Phi)
    \sim_{\mathrm{ret}}
    (Q,\widehat\Psi),
\]
if they are connected by a finite zigzag of strict state retractions.

Equivalently,
\(\sim_{\mathrm{ret}}\) is the connected-component relation in the
hom-category
\[
    \mathbf{RetPres}^{\mathrm{cen}}_{\mathsf K}
    (\mathbb A,\mathbb B).
\]
\end{definition}

\begin{proposition}[Isolation of the empty retractive class]
\label{prop:empty-retractive-class-isolated}
Suppose that the initial object \(0\) of \(\mathcal E\) is strict.  If a
strict state retraction has either the empty presentation as its domain or
the empty presentation as its codomain, then its other endpoint also has
initial carrier.

Consequently, a finite zigzag of strict state retractions connects an empty
presentation only to presentations with initial carrier.
\end{proposition}

\begin{proof}
Suppose first that
\[
    h:P\longrightarrow0
\]
is a strict state retraction.  Its underlying carrier map is a morphism
\(P\to0\).  Since \(0\) is strict, this map is an isomorphism and
\(P\cong0\).

Suppose instead that
\[
    h:0\longrightarrow P
\]
is a strict state retraction.  A strict section of \(h\) is in particular a
carrier map
\[
    P\longrightarrow0.
\]
Strictness of \(0\) again gives \(P\cong0\).

The zigzag statement follows by induction on its length.
\end{proof}

\begin{corollary}[Noncollapse in a feedback span base]
\label{cor:retractive-quotient-noncollapse}
Let \(\mathcal E\) be a feedback span base.  Then the initial object is
strict.  Hence the empty presentation cannot be retractively
presentation-equivalent to a presentation with noninitial carrier.

In particular, if \(A_0\not\cong0\), then the empty endopresentation of
\(\mathbb A\) and the identity presentation of \(\mathbb A\) determine
distinct classes.
\end{corollary}

\begin{proof}
The initial object is the empty coproduct.  Pullback stability of the empty
coproduct says that for every map \(X\to0\), the canonical map
\[
    0\longrightarrow X
\]
obtained by pulling back the empty coproduct decomposition is an isomorphism.
Thus \(0\) is strict.

The remaining statements follow from
Proposition~\ref{prop:empty-retractive-class-isolated}.  The identity
presentation has carrier \(A_0\), so it cannot lie in the empty class when
\(A_0\) is noninitial.
\end{proof}

\begin{theorem}[Retractive state-presentation quotient]
\label{thm:retractive-state-presentation-quotient}
Retractive presentation equivalence is a congruence for identities and
pipeline composition.  There is therefore a category
\[
    \mathbf{Mly}^{\mathrm{cen},\mathrm{ret}}_{\mathsf K}(\mathcal E)
\]
whose objects are internal categories and whose morphisms are
\(\sim_{\mathrm{ret}}\)-classes of central effectful Mealy morphisms.
\end{theorem}

\begin{proof}
Horizontal composition in
\(\mathbf{RetPres}^{\mathrm{cen}}_{\mathsf K}(\mathcal E)\)
carries strict state retractions to strict state retractions and therefore
carries zigzags to zigzags.

The canonical associators and unitors are strict state isomorphisms and hence
strict state retractions.  Their endpoints therefore belong to the same
connected component.  Consequently composition on connected components is
strictly associative and unital.
\end{proof}

\section{Retract-invariant state hiding}
\label{sec:retract-invariant-state-hiding}

Let
\[
    \operatorname{Cat}_1(\mathcal E)
\]
denote the ordinary category whose objects are internal categories in
\(\mathcal E\) and whose morphisms are internal functors.  Natural
transformations are not retained at this categorical level.

Every internal functor
\[
    F:\mathbb A\longrightarrow\mathbb B
\]
determines a pure stateless effectful Mealy morphism, written
\[
    J^{\mathrm{st}}_{\mathsf K}(F):
    \mathbb A\rightsquigarrow_{\mathsf K}\mathbb B.
\]
Its carrier is the map-representable span \(F_0!\), and its arrow map is
inserted by \(\eta^{\mathsf K}\).

\begin{definition}[Typed retract-invariant state-hiding interpretation]
\label{def:retract-invariant-state-hiding-interpretation}
Let \(\mathcal D\) be a category.  A
\textbf{typed retract-invariant state-hiding interpretation} of central
effectful Mealy presentations in \(\mathcal D\) consists of:
\begin{enumerate}
    \item a functor
    \[
        J:\operatorname{Cat}_1(\mathcal E)\longrightarrow\mathcal D;
    \]

    \item for every central effectful Mealy presentation
    \[
        (P,\widehat\Phi):
        \mathbb A\rightsquigarrow_{\mathsf K}\mathbb B,
    \]
    a morphism
    \[
        \operatorname{hide}_P(\widehat\Phi):
        J\mathbb A\longrightarrow J\mathbb B.
    \]
\end{enumerate}
These data satisfy the following equations.

\begin{enumerate}
    \item \textbf{Pure presentations.}
    For every internal functor \(F:\mathbb A\to\mathbb B\),
    \[
        \operatorname{hide}_{F_0!}
        \bigl(
            \widehat\Phi_{J^{\mathrm{st}}_{\mathsf K}(F)}
        \bigr)
        =
        J(F).
    \]

    \item \textbf{Tightening.}
    Let
    \[
        L:\mathbb A'\to\mathbb A,
        \qquad
        R:\mathbb B\to\mathbb B'
    \]
    be internal functors.  Then
    \[
\begin{aligned}
&
\operatorname{hide}_{R_0!\circ P\circ L_0!}
\Bigl(
    \widehat\Phi_{
        J^{\mathrm{st}}_{\mathsf K}(R)
        \circ
        (P,\widehat\Phi)
        \circ
        J^{\mathrm{st}}_{\mathsf K}(L)
    }
\Bigr)
\\
&\qquad=
J(R)
\circ
\operatorname{hide}_P(\widehat\Phi)
\circ
J(L).
\end{aligned}
    \]

    \item \textbf{Joining.}
    For composable presentations
    \[
        (P,\widehat\Phi_P),
        \qquad
        (Q,\widehat\Phi_Q),
    \]
    \[
        \operatorname{hide}_{Q\circ P}
        (\widehat\Phi_{Q\circ P})
        =
        \operatorname{hide}_Q(\widehat\Phi_Q)
        \circ
        \operatorname{hide}_P(\widehat\Phi_P).
    \]

    \item \textbf{Retract invariance.}
    If
    \[
        h:
        (P,\widehat\Phi)
        \longrightarrow
        (Q,\widehat\Psi)
    \]
    is a strict state retraction, then
    \[
        \operatorname{hide}_P(\widehat\Phi)
        =
        \operatorname{hide}_Q(\widehat\Psi).
    \]
\end{enumerate}
\end{definition}

The joining equation hides two successive state carriers as their composite
span.  Tightening transports a hidden presentation along external internal
functors.  Retract invariance says that the interpreted external process is
unchanged when a presentation is replaced by a strict retract.

The tightening equation follows from the pure-presentation and joining
equations, but it is retained explicitly because it records the intended
typed covariance of the state-hiding interpretation.

\begin{proposition}[Canonical retract-invariant state hiding]
\label{prop:canonical-retract-invariant-state-hiding}
The quotient category
\[
    \mathbf{Mly}^{\mathrm{cen},\mathrm{ret}}_{\mathsf K}(\mathcal E)
\]
carries a canonical typed retract-invariant state-hiding interpretation.  The
functor
\[
    J^{\mathrm{st}}_{\mathsf K}:
    \operatorname{Cat}_1(\mathcal E)
    \longrightarrow
    \mathbf{Mly}^{\mathrm{cen},\mathrm{ret}}_{\mathsf K}(\mathcal E)
\]
sends an internal functor to the class of its pure stateless effectful Mealy
morphism, and
\[
    \operatorname{hide}_P(\widehat\Phi)
    :=
    [(P,\widehat\Phi)]_{\mathrm{ret}}.
\]
\end{proposition}

\begin{proof}
Pure stateless presentations compose as internal functors.  Tightening and
joining are the corresponding composition equations in the quotient.
Retract invariance holds because the endpoints of a strict state retraction
lie in the same connected component of the retractive-presentation
hom-category.
\end{proof}

\begin{theorem}[Universal retractive state-presentation quotient]
\label{thm:universal-retractive-state-presentation-quotient}
Let \(\mathcal D\) carry a typed retract-invariant state-hiding
interpretation \((J,\operatorname{hide})\).  There is a unique functor
\[
    \overline J:
    \mathbf{Mly}^{\mathrm{cen},\mathrm{ret}}_{\mathsf K}(\mathcal E)
    \longrightarrow
    \mathcal D
\]
such that
\[
    \overline J\circ J^{\mathrm{st}}_{\mathsf K}
    =
    J
\]
and
\[
    \overline J
    \bigl(
        [(P,\widehat\Phi)]_{\mathrm{ret}}
    \bigr)
    =
    \operatorname{hide}_P(\widehat\Phi).
\]
\end{theorem}

\begin{proof}
Define
\[
    \overline J(\mathbb A)
    :=
    J\mathbb A
\]
on objects and
\[
    \overline J
    \bigl(
        [(P,\widehat\Phi)]_{\mathrm{ret}}
    \bigr)
    :=
    \operatorname{hide}_P(\widehat\Phi)
\]
on morphisms.

Retract invariance identifies the endpoints of every strict state retraction
and therefore every finite zigzag of such retractions, so the definition is
well-defined.

The pure-presentation equation gives
\[
    \overline J
    \bigl(
        J^{\mathrm{st}}_{\mathsf K}(F)
    \bigr)
    =
    J(F).
\]
The joining equation gives
\[
\begin{aligned}
\overline J
\bigl(
    [(Q,\widehat\Phi_Q)]_{\mathrm{ret}}
    \circ
    [(P,\widehat\Phi_P)]_{\mathrm{ret}}
\bigr)
&=
\operatorname{hide}_{Q\circ P}
(\widehat\Phi_{Q\circ P})
\\
&=
\operatorname{hide}_Q(\widehat\Phi_Q)
\circ
\operatorname{hide}_P(\widehat\Phi_P).
\end{aligned}
\]
The identity equation is the pure-presentation equation for the identity
internal functor.  Thus \(\overline J\) is a functor.

Every morphism in the quotient is represented by a presentation.  Hence any
functor preserving the state-hiding interpretation must have the displayed
value on every morphism, proving uniqueness.
\end{proof}

\begin{corollary}[Retractive presentation congruence]
\label{cor:retractive-presentation-congruence}
Retractive presentation equivalence is the smallest composition congruence on
central effectful Mealy morphisms that identifies the endpoints of every
strict state retraction.
\end{corollary}

\begin{proof}
The relation is generated by finite zigzags of strict state retractions and
is a composition congruence by
Theorem~\ref{thm:retractive-state-presentation-quotient}.  Every composition
congruence identifying the endpoints of each strict state retraction
therefore identifies the endpoints of every generating zigzag.
\end{proof}

\begin{proposition}[Unpointed presentations do not determine pointed feedback]
\label{prop:unpointed-not-pointed-feedback}
The unpointed state presentations used in this chapter do not determine the
pointed uniform-feedback operation of
\cite{BonchiDiLavoreRomanEffectfulMealyMachines}, even in
\(\mathbf{Set}\).
\end{proposition}

\begin{proof}
Let
\[
    S=\{0,1\},
    \qquad
    X=1,
    \qquad
    Y=\{0,1\},
\]
and define
\[
    f:S\times X\longrightarrow S\times Y
\]
by
\[
    f(s,*)=(s,s).
\]
With initial state \(0\), pointed feedback produces the constant output
\(0\).  With initial state \(1\), it produces the constant output \(1\).
The underlying unpointed state object and transition map are identical.
Therefore the unpointed presentation does not determine the pointed
feedback process.
\end{proof}

\section{Effectful finite-path feedback}
\label{sec:effectful-finite-path-feedback}

Assume throughout this section that \(\mathcal E\) is a feedback span base in
the sense of Definition~9.2.1.  Thus the finite and countable coproducts used
below are disjoint and pullback-stable, and the required countable
reindexing isomorphisms are available.

We distinguish:
\[
    +
\]
for the coproduct tensor, which changes source and target interfaces, from
\[
    \boxplus
\]
for the homwise coproduct of parallel spans, machines, cells, and
presentations.

\subsection{Homwise sums of effectful span cells}
\label{subsec:homwise-sums-effectful-cells}

Let \((P_i)_{i\in I}\) and \((Q_i)_{i\in I}\) be finite or countable families
of parallel spans from \(X\) to \(Y\), and let
\[
    \widehat f_i:
    P_i\Rrightarrow_{\mathsf K}Q_i
\]
be effectful cells.  Define
\[
    \boxplus_{i\in I}\widehat f_i:
    \boxplus_{i\in I}P_i
    \Rrightarrow_{\mathsf K}
    \boxplus_{i\in I}Q_i
\]
by the copairing whose \(i\)-th component is
\[
    P_i
    \xrightarrow{\widehat f_i}
    \mathsf K(Q_i)
    \xrightarrow{\mathsf K(\iota_i)}
    \mathsf K
    \left(
        \coprod_{j\in I}Q_j
    \right).
\]

\begin{proposition}[Central homwise sums of span cells]
\label{prop:central-homwise-sums}
If every \(\widehat f_i\) is central, then
\[
    \boxplus_i\widehat f_i
\]
is central.

For aligned composable families,
\[
    \left(
        \boxplus_i\widehat g_i
    \right)
    \bullet
    \left(
        \boxplus_i\widehat f_i
    \right)
    =
    \boxplus_i
    \left(
        \widehat g_i\bullet\widehat f_i
    \right).
\]

Horizontal composition is bilinear:
\[
    \widehat g
    \ast
    \left(
        \boxplus_i\widehat f_i
    \right)
    \cong
    \boxplus_i
    \left(
        \widehat g\ast\widehat f_i
    \right),
\]
\[
    \left(
        \boxplus_j\widehat g_j
    \right)
    \ast
    \widehat f
    \cong
    \boxplus_j
    \left(
        \widehat g_j\ast\widehat f
    \right).
\]
Consequently,
\[
    \left(
        \boxplus_j\widehat g_j
    \right)
    \ast
    \left(
        \boxplus_i\widehat f_i
    \right)
    \cong
    \boxplus_{i,j}
    \left(
        \widehat g_j\ast\widehat f_i
    \right).
\]
\end{proposition}

\begin{proof}
After arbitrary reindexing, pullback stability gives
\[
    \left(
        \coprod_iP_i
    \right)
    \times R
    \cong
    \coprod_i(P_i\times R).
\]
Restricting the two centrality composites along the \(i\)-th coproduct
injection gives the centrality equation for \(\widehat f_i\), followed by the
same pure target injection.  Since the coproduct injections are jointly
epimorphic, the two composites agree.

Compatibility with vertical composition follows by restricting to each
summand, using naturality of \(\mu\), and applying the Kleisli composition
equation there.

For horizontal composition, pullback distributes over the coproducts.
Hence the composite carrier decomposes into the indicated summands.  On each
summand, the transition cell is the corresponding horizontal composite.
The canonical coproduct reindexing gives the displayed isomorphisms.
\end{proof}

\subsection{Homwise sums of effectful Mealy presentations}
\label{subsec:homwise-sums-effectful-mealy}

Let
\[
    (P_i,\widehat\Phi_i):
    \mathbb A\rightsquigarrow_{\mathsf K}\mathbb B
\]
be a finite or countable family of parallel central effectful Mealy
morphisms.  Put
\[
    P:=\coprod_iP_i
\]
with the induced span legs.  Pullback stability gives canonical
isomorphisms
\[
    P\circ A
    \cong
    \coprod_i(P_i\circ A),
    \qquad
    B\circ P
    \cong
    \coprod_i(B\circ P_i).
\]

Define
\[
    \boxplus_i\widehat\Phi_i:
    P\circ A
    \longrightarrow
    \mathsf K(B\circ P)
\]
by the copairing whose \(i\)-th restriction is
\[
    P_i\circ A
    \xrightarrow{\widehat\Phi_i}
    \mathsf K(B\circ P_i)
    \xrightarrow{\mathsf K(B\circ\iota_i)}
    \mathsf K(B\circ P).
\]

\begin{proposition}[Homwise sums of effectful Mealy morphisms]
\label{prop:homwise-sums-effectful-mealy}
The pair
\[
    \boxplus_i(P_i,\widehat\Phi_i)
    :=
    \left(
        \coprod_iP_i,
        \boxplus_i\widehat\Phi_i
    \right)
\]
is a central effectful Mealy morphism
\[
    \mathbb A\rightsquigarrow_{\mathsf K}\mathbb B.
\]

Let
\[
    \widehat\Theta_i:
    (P_i,\widehat\Phi_i)
    \Rightarrow
    (Q_i,\widehat\Psi_i)
\]
be a family of central effectful simulations.  Their homwise sum is the map
\[
    \boxplus_i\widehat\Theta_i:
    \coprod_iQ_i
    \longrightarrow
    \mathsf K
    \left(
        B\circ\coprod_iP_i
    \right)
\]
whose \(i\)-th restriction is
\[
    Q_i
    \xrightarrow{\widehat\Theta_i}
    \mathsf K(B\circ P_i)
    \xrightarrow{\mathsf K(B\circ\iota_i)}
    \mathsf K
    \left(
        B\circ\coprod_jP_j
    \right).
\]
This is a central effectful simulation
\[
    \boxplus_i(P_i,\widehat\Phi_i)
    \Rightarrow
    \boxplus_i(Q_i,\widehat\Psi_i).
\]

For vertically composable aligned families,
\[
    \left(
        \boxplus_i\widehat\Lambda_i
    \right)
    \circ_v
    \left(
        \boxplus_i\widehat\Theta_i
    \right)
    =
    \boxplus_i
    \left(
        \widehat\Lambda_i\circ_v\widehat\Theta_i
    \right).
\]
Identity simulations are preserved.
\end{proposition}

\begin{proof}
The effectful Mealy unit equation for
\(\boxplus_i\widehat\Phi_i\) can be checked after restriction to every
summand \(P_i\circ A\).  On that summand it is the unit equation for
\(\widehat\Phi_i\), followed by the corresponding pure coproduct injection.
Joint epimorphy of the coproduct injections gives the global equation.

The multiplication equation is checked similarly.  Pullback stability
decomposes the iterated input-transition object into the summands belonging
to the \(P_i\).  Restriction to the \(i\)-th summand gives the multiplication
equation for \(\widehat\Phi_i\).  No cross-summand state transition occurs,
because the carrier injections preserve both interface legs.

Centrality follows from
Proposition~\ref{prop:central-homwise-sums}.

For simulations, restrict the simulation equation for
\(\boxplus_i\widehat\Theta_i\) to \(Q_i\circ A\).  Both sides reduce to the
simulation equation for \(\widehat\Theta_i\), followed by the same pure
injection into the total target carrier.  Hence the equation holds globally.
Centrality again follows from
Proposition~\ref{prop:central-homwise-sums}.

The vertical-composition equation follows by restriction to each summand and
Proposition~\ref{prop:central-homwise-sums}.  Identity simulations are
copairings of the summandwise identities.
\end{proof}

\begin{proposition}[Pipeline distributivity over homwise sums]
\label{prop:pipeline-distributes-homwise-sums}
Let
\[
    P_i:
    \mathbb A\rightsquigarrow_{\mathsf K}\mathbb B,
    \qquad
    Q:
    \mathbb B\rightsquigarrow_{\mathsf K}\mathbb C
\]
be central effectful Mealy morphisms.  There is a canonical strict state
isomorphism
\[
    Q\circ\left(\boxplus_iP_i\right)
    \xrightarrow{\cong}
    \boxplus_i(Q\circ P_i).
\]
Similarly, for
\[
    P:
    \mathbb A\rightsquigarrow_{\mathsf K}\mathbb B,
    \qquad
    Q_j:
    \mathbb B\rightsquigarrow_{\mathsf K}\mathbb C,
\]
there is a canonical strict state isomorphism
\[
    \left(\boxplus_jQ_j\right)\circ P
    \xrightarrow{\cong}
    \boxplus_j(Q_j\circ P).
\]
Consequently,
\[
    \left(\boxplus_jQ_j\right)
    \circ
    \left(\boxplus_iP_i\right)
    \xrightarrow{\cong}
    \boxplus_{i,j}(Q_j\circ P_i).
\]

The same comparisons are natural with respect to central simulations and
strict state homomorphisms.
\end{proposition}

\begin{proof}
For the first comparison, pullback stability gives
\[
    \left(
        \coprod_iP_i
    \right)
    \times_{B_0}Q
    \cong
    \coprod_i
    \left(
        P_i\times_{B_0}Q
    \right).
\]
Under this carrier isomorphism, the pipeline transition restricts on the
\(i\)-th summand to the pipeline transition of \(Q\circ P_i\).  Hence the
carrier isomorphism satisfies the strict transition equation.

The second comparison is identical, and the double-family comparison is
their composite.  Naturality follows from naturality of the
pullback--coproduct comparison and the summandwise definitions of simulation
and strict-map composition.
\end{proof}

\subsection{The coproduct tensor}
\label{subsec:effectful-coproduct-tensor}

Let
\[
    \Phi:
    \mathbb A\rightsquigarrow_{\mathsf K}\mathbb B,
    \qquad
    \Psi:
    \mathbb C\rightsquigarrow_{\mathsf K}\mathbb D
\]
be central effectful Mealy morphisms with carriers \(P\) and \(Q\).

There are canonical block decompositions
\[
    (P+Q)\circ(A+C)
    \cong
    (P\circ A)+(Q\circ C),
\]
\[
    (B+D)\circ(P+Q)
    \cong
    (B\circ P)+(D\circ Q).
\]
The mixed terms are initial because the coproduct internal categories have no
cross-summand arrows.

Let
\[
    j_{AB}:
    A_0\times B_0
    \longrightarrow
    (A_0+C_0)\times(B_0+D_0)
\]
and
\[
    j_{CD}:
    C_0\times D_0
    \longrightarrow
    (A_0+C_0)\times(B_0+D_0)
\]
be the diagonal block inclusions.

\begin{definition}[Coproduct tensor of effectful Mealy morphisms]
\label{def:coproduct-tensor-effectful-mealy}
The coproduct tensor
\[
    \Phi+\Psi:
    \mathbb A+\mathbb C
    \rightsquigarrow_{\mathsf K}
    \mathbb B+\mathbb D
\]
has carrier \(P+Q\).  Its transition map is the copairing of the two
dependent-sum transports
\[
    \Sigma_{j_{AB}}^{\mathsf K}\widehat\Phi,
    \qquad
    \Sigma_{j_{CD}}^{\mathsf K}\widehat\Psi
\]
under the canonical block decompositions above.
\end{definition}

\begin{theorem}[Coproduct symmetric monoidal effectful Mealy bicategory]
\label{thm:central-effectful-coproduct-monoidal-bicategory}
There is a symmetric monoidal bicategory
\[
    \left(
        \mathbf{Mly}^{\mathrm{cen},+}_{\mathsf K,\omega}(\mathcal E),
        +,
        0
    \right)
\]
whose underlying bicategory is
\(\mathbf{Mly}^{\mathrm{cen}}_{\mathsf K}(\mathcal E)\) and whose tensor is
the chosen coproduct tensor.
\end{theorem}

\begin{proof}
Dependent-sum transport preserves centrality, and homwise coproduct preserves
centrality.  Therefore the transition of
\(\Phi+\Psi\) is central.

The Mealy unit and multiplication equations can be checked after restriction
to the two diagonal blocks.  They then reduce respectively to the equations
for \(\Phi\) and \(\Psi\).  The mixed blocks are initial.

For pipeline composition, pullback stability gives
\[
    (P+Q)\times_{B_0+D_0}(P'+Q')
    \cong
    (P\times_{B_0}P')
    +
    (Q\times_{D_0}Q').
\]
Under this isomorphism, the two pipeline transitions agree on both summands.
Thus coproduct is pseudofunctorial.

The tensor of simulations is defined by the same blockwise copairing.
Vertical and horizontal functoriality follow from
Propositions~\ref{prop:central-homwise-sums} and
\ref{prop:homwise-sums-effectful-mealy}.  The coproduct associator, unitors,
and symmetry are pure, and their coherence reduces to coproduct coherence in
\(\mathcal E\).
\end{proof}

Passing to isomorphism classes in the hom-categories gives a symmetric
monoidal category, written
\[
    \left(
        \operatorname{Mly}^{\mathrm{cen},+}_{\mathsf K,\omega,0}
        (\mathcal E),
        +,
        0
    \right).
\]

\subsection{Effectful path objects}
\label{subsec:effectful-path-objects}

Let
\[
    U:
    \mathbb C\rightsquigarrow_{\mathsf K}\mathbb C
\]
be a central effectful Mealy endomorphism.  Define
\[
    U^{[0]}:=1_{\mathbb C},
    \qquad
    U^{[n+1]}:=U^{[n]}\circ U.
\]

\begin{definition}[Effectful finite-path object]
\label{def:effectful-finite-path-object}
The \textbf{effectful finite-path Mealy morphism} generated by \(U\) is
\[
    \operatorname{Path}_{\mathsf K}(U)
    :=
    \boxplus_{n\geq0}U^{[n]}.
\]
Its canonical generator is the one-length coproduct injection
\[
    \iota_1:
    U
    \longrightarrow
    \operatorname{Path}_{\mathsf K}(U).
\]
\end{definition}

Its carrier is the coproduct of the finite composite carriers.  On the
\(n\)-th summand, its transition is the \(n\)-fold central pipeline
composite.  Path witnesses remain explicit in the carrier, while the effects
generated along a fixed path are combined by the multiplication of
\(\mathsf K\).

\begin{lemma}[Homwise coproducts of strict presentations]
\label{lem:strict-presentation-homwise-coproducts}
For a finite or countable family
\[
    (U_i)_{i\in I}
\]
of parallel central effectful Mealy morphisms, the homwise sum
\[
    \boxplus_{i\in I}U_i
\]
is their coproduct in
\[
    \operatorname{Pres}^{\mathrm{cen}}_{\mathsf K}
    (\mathbb A,\mathbb B).
\]
The coproduct injection
\[
    U_i\longrightarrow\boxplus_jU_j
\]
is a strict state homomorphism.

A carrier map
\[
    h:
    \boxplus_iU_i
    \longrightarrow
    V
\]
is a strict state homomorphism if and only if every restriction
\[
    h_i:
    U_i
    \longrightarrow
    V
\]
is a strict state homomorphism.
\end{lemma}

\begin{proof}
A carrier map from the coproduct carrier is equivalently a family of carrier
maps from its summands.  Pullback stability of the coproduct decomposes the
transition object of \(\boxplus_iU_i\) into the corresponding transition
objects of the \(U_i\).

The transition equation for the \(i\)-th coproduct injection restricts to the
identity transition equation on \(U_i\), so each injection is strict.

Restricting the strict transition equation for \(h\) along the \(i\)-th
coproduct injection gives the strict transition equation for \(h_i\).
Conversely, if every \(h_i\) is strict, the two sides of the strict transition
equation for \(h\) agree after restriction to every summand.  Since the
coproduct injections are jointly epimorphic, the two sides agree globally.
The ordinary coproduct universal property of the carrier then gives the
required coproduct in the strict-presentation hom-category.
\end{proof}

\begin{proposition}[Effectful path multiplication]
\label{prop:effectful-path-multiplication}
Finite-length concatenation induces a strict state homomorphism
\[
    \operatorname{mult}_U:
    \operatorname{Path}_{\mathsf K}(U)
    \circ
    \operatorname{Path}_{\mathsf K}(U)
    \longrightarrow
    \operatorname{Path}_{\mathsf K}(U)
\]
whose restriction to the \((m,n)\)-summand is
\[
    U^{[m]}\circ U^{[n]}
    \xrightarrow{\cong}
    U^{[m+n]}
    \xrightarrow{\iota_{m+n}}
    \operatorname{Path}_{\mathsf K}(U).
\]
Together with the zero-length inclusion
\[
    \operatorname{unit}_U:
    1_{\mathbb C}
    \longrightarrow
    \operatorname{Path}_{\mathsf K}(U),
\]
this makes
\(\operatorname{Path}_{\mathsf K}(U)\) a monoid object in
\[
    \operatorname{Pres}^{\mathrm{cen}}_{\mathsf K}
    (\mathbb C,\mathbb C)
\]
under pipeline composition.
\end{proposition}

\begin{proof}
By Proposition~\ref{prop:pipeline-distributes-homwise-sums},
\[
\operatorname{Path}_{\mathsf K}(U)
\circ
\operatorname{Path}_{\mathsf K}(U)
\cong
\boxplus_{m,n\geq0}
\left(
    U^{[m]}\circ U^{[n]}
\right).
\]
The canonical associators identify each summand with \(U^{[m+n]}\).
Both sides of the strict transition equation on this summand are the same
\((m+n)\)-fold ordered pipeline transition, differing only by the canonical
pullback associator.  Thus each summand map is strict, and
Lemma~\ref{lem:strict-presentation-homwise-coproducts} makes their copairing
strict.

The monoid equations are read using the canonical pipeline associator and
unitors.  On every length summand, associativity reduces to
\[
    (m+n)+r=m+(n+r)
\]
and the coherence of the pipeline associator.  The zero-length summand is the
pipeline unit.
\end{proof}

\begin{theorem}[Free effectful path monoid]
\label{thm:free-effectful-path-monoid}
The monoid
\[
    \operatorname{Path}_{\mathsf K}(U)
\]
with generator
\[
    \iota_1:
    U\to\operatorname{Path}_{\mathsf K}(U)
\]
is the free monoid on \(U\) in
\[
    \operatorname{Pres}^{\mathrm{cen}}_{\mathsf K}
    (\mathbb C,\mathbb C).
\]
Explicitly, for every monoid object \(M\) in this endomorphism category,
restriction along \(\iota_1\) induces a natural bijection
\[
\operatorname{Mon}
\left(
    \operatorname{Path}_{\mathsf K}(U),
    M
\right)
\cong
\operatorname{Pres}^{\mathrm{cen}}_{\mathsf K}
(\mathbb C,\mathbb C)(U,M).
\]
\end{theorem}

\begin{proof}
Let \(M\) have multiplication \(\mu_M\) and unit \(\eta_M\), and let
\[
    f:U\to M
\]
be a strict state homomorphism.  Define
\[
    f_0:=\eta_M
\]
and recursively
\[
    f_{n+1}
    :=
    \mu_M
    \circ
    (f_n\circ f):
    U^{[n+1]}\to M.
\]
Every \(f_n\) is strict by closure under pipeline composition.  By
Lemma~\ref{lem:strict-presentation-homwise-coproducts}, their copairing gives
a unique strict state homomorphism
\[
    f^\sharp:
    \operatorname{Path}_{\mathsf K}(U)
    \longrightarrow
    M.
\]
The recursive definition gives the monoid equations on each length summand,
using
Proposition~\ref{prop:pipeline-distributes-homwise-sums}
and the canonical pipeline associators and unitors.  Hence
\(f^\sharp\) is a monoid morphism, and its restriction to the one-length
summand is \(f\).

Conversely, let
\[
    g:
    \operatorname{Path}_{\mathsf K}(U)
    \longrightarrow
    M
\]
be a monoid morphism whose restriction to the one-length summand is \(f\).
The unit equation forces its zero-length component to be \(f_0\), and the
multiplication equation forces its \((n+1)\)-length component to be
\[
    \mu_M\circ(f_n\circ f).
\]
Induction gives equality with \(f_n\) on every summand.  The coproduct
universal property from
Lemma~\ref{lem:strict-presentation-homwise-coproducts} therefore gives
\(g=f^\sharp\).

The two assignments are inverse and natural in \(M\).
\end{proof}

Under the contravariant displayed-simulation convention, the same carrier
maps appear in the opposite direction.  The monoid statement therefore
belongs naturally to the covariant strict-presentation category.

\subsection{Block decomposition and trace}
\label{subsec:effectful-block-trace}

Let
\[
    \Phi:
    \mathbb A+\mathbb C
    \rightsquigarrow_{\mathsf K}
    \mathbb B+\mathbb C
\]
be a central effectful Mealy morphism.  Pullback along the source and target
coproduct injections produces four blocks
\[
    \Phi_{AB},
    \qquad
    \Phi_{AC},
    \qquad
    \Phi_{CB},
    \qquad
    \Phi_{CC}.
\]
The first index records the source summand, and the second records the target
summand.

\begin{proposition}[Effectful block restriction]
\label{prop:effectful-block-restriction}
Each of the four blocks is a central effectful Mealy morphism of the indicated
type.

Block restriction also sends:
\begin{enumerate}
    \item a central effectful simulation to central simulations between the
    corresponding blocks;

    \item a strict state homomorphism to strict state homomorphisms between
    the corresponding blocks;

    \item a strict state retraction, together with any chosen strict section,
    to blockwise strict state retractions and blockwise strict sections.
\end{enumerate}
\end{proposition}

\begin{proof}
Pull back the carrier and transition map along the corresponding source and
target block inclusions.  Fibredness gives
\[
    j^*\mathsf K(Q)
    \cong
    \mathsf K(j^*Q).
\]
Because the coproduct internal categories contain no cross-summand arrows,
an output-comparison arrow whose target lies in one summand has its source in
the same summand.  Therefore the transition restricts to the selected block.
The Mealy equations are preserved by pullback.  Centrality is preserved
because central arrows are stable under reindexing.

The same pullback construction restricts a simulation equation to each
block, giving a block simulation.  A strict carrier map preserves both
interface legs, so it maps each carrier block into the corresponding block.
Pulling back its strict transition equation gives the blockwise strict
transition equation.

Finally, if
\[
    h\circ s=1,
\]
then pulling this equality back to any block gives
\[
    h_{XY}\circ s_{XY}=1.
\]
Thus block restriction preserves chosen retraction data.
\end{proof}

\begin{definition}[Effectful coproduct trace]
\label{def:effectful-coproduct-trace}
Define
\[
    \operatorname{Tr}^{\mathbb C}_{\mathsf K}(\Phi)
    :=
    \Phi_{AB}
    \boxplus
    \Phi_{CB}
    \circ
    \operatorname{Path}_{\mathsf K}(\Phi_{CC})
    \circ
    \Phi_{AC}.
\]
Equivalently,
\[
    \operatorname{Tr}^{\mathbb C}_{\mathsf K}(\Phi)
    =
    \Phi_{AB}
    \boxplus
    \boxplus_{n\geq0}
    \left(
        \Phi_{CB}
        \circ
        \Phi_{CC}^{[n]}
        \circ
        \Phi_{AC}
    \right).
\]
\end{definition}

The first summand is direct visible execution.  The \(n\)-th hidden summand
enters the hidden interface, performs \(n\) hidden steps, and exits.  Effects
are sequenced along the hidden path in their edge order.

\begin{lemma}[Decorated hidden-path normal form]
\label{lem:decorated-hidden-path-normal-form}
Every expression built from:
\begin{enumerate}
    \item pipeline composition;

    \item homwise coproduct;

    \item coproduct tensor;

    \item coproduct symmetries;

    \item finitely many applications of
    \(\operatorname{Tr}^{(-)}_{\mathsf K}\);
\end{enumerate}
admits a canonical expansion as a homwise coproduct indexed by finite
hidden-interface paths.

For each path, the corresponding effectful transition is the ordered
pipeline composite of the block transitions decorating its edges.
Canonical reindexings of hidden paths induce pure strict isomorphisms between
the corresponding effectful summands.
\end{lemma}

\begin{proof}
Proceed by structural induction.

Pipeline composition concatenates path expressions.
Proposition~\ref{prop:pipeline-distributes-homwise-sums} yields the coproduct
indexed by pairs of paths, and the canonical concatenation reindexing sends
each pair to its concatenated path.  The transition on that summand is the
ordered pipeline composite.

Homwise coproduct forms the disjoint union of the path families without
changing their external interfaces.  Coproduct tensor forms the disjoint
union of path families in separate interface blocks.  Coproduct symmetry
renames visible and hidden summands and preserves the edge order.

An application of trace replaces the hidden endomorphism block by
\[
    \boxplus_{n\geq0}\Phi_{CC}^{[n]},
\]
so it inserts all finite hidden paths in that block.

The required countable reindexing isomorphisms exist by the feedback-base
conditions.  On corresponding summands, the transition maps are the same
ordered list of central pipeline transitions, up to the canonical pullback
associators and unitors.  These coherence maps are pure strict
isomorphisms.
\end{proof}

\begin{theorem}[Effectful traced Mealy category]
\label{thm:effectful-traced-mealy-category}
The symmetric monoidal category
\[
    \left(
        \operatorname{Mly}^{\mathrm{cen},+}_{\mathsf K,\omega,0}
        (\mathcal E),
        +,
        0
    \right)
\]
is traced.  Its trace is
\[
    \operatorname{Tr}^{\mathbb C}_{\mathsf K}
\]
from Definition~\ref{def:effectful-coproduct-trace}.
\end{theorem}

\begin{proof}
We verify the Joyal--Street--Verity trace equations by decorated hidden-path
normal forms.

\textbf{Naturality.}
Precomposition and postcomposition by visible morphisms affect only the
first and last visible edges of every path.  Bilinearity of composition over
homwise coproducts gives
\[
    g\circ
    \operatorname{Tr}^{\mathbb C}_{\mathsf K}(\Phi)
    \circ f
    =
    \operatorname{Tr}^{\mathbb C}_{\mathsf K}
    \bigl(
        (g+1_{\mathbb C})
        \circ
        \Phi
        \circ
        (f+1_{\mathbb C})
    \bigr).
\]

\textbf{Dinaturality.}
Let
\[
    \Phi:
    \mathbb A+\mathbb C
    \rightsquigarrow_{\mathsf K}
    \mathbb B+\mathbb D
\]
and
\[
    g:\mathbb D\rightsquigarrow_{\mathsf K}\mathbb C.
\]
A hidden path on the left-hand side
\[
    \operatorname{Tr}^{\mathbb C}
    \bigl(
        (1_{\mathbb B}+g)\circ\Phi
    \bigr)
\]
has ordered edge sequence
\[
    \Phi_{AD},
    \ g,
    \ \Phi_{CD},
    \ g,
    \ \ldots,
    \ \Phi_{CB}.
\]
The corresponding path on the right-hand side
\[
    \operatorname{Tr}^{\mathbb D}
    \bigl(
        \Phi\circ(1_{\mathbb A}+g)
    \bigr)
\]
has exactly the same ordered sequence, with the cuts between hidden steps
placed differently.  The two path families are canonically bijective, and
Lemma~\ref{lem:decorated-hidden-path-normal-form} identifies the decorated
summands.

\textbf{Vanishing over \(0\).}
There are no nontrivial hidden \(0\)-paths.  The trace formula reduces to the
visible block, giving
\[
    \operatorname{Tr}^{0}_{\mathsf K}(\Phi)
    =
    \Phi.
\]

\textbf{Double vanishing.}
A finite path through
\[
    \mathbb C+\mathbb D
\]
is equivalently a finite list of maximal \(\mathbb C\)-segments and
\(\mathbb D\)-segments.  The feedback-base reindexing by finite lists gives
the canonical bijection between
\[
    \operatorname{Tr}^{\mathbb C+\mathbb D}_{\mathsf K}(\Phi)
\]
and the iterated trace
\[
    \operatorname{Tr}^{\mathbb C}_{\mathsf K}
    \left(
        \operatorname{Tr}^{\mathbb D}_{\mathsf K}(\Phi)
    \right).
\]
The reindexing preserves the ordered edge sequence, so corresponding
effectful transitions agree.

\textbf{Superposing.}
A path in
\[
    \Phi+\Psi
\]
lies entirely in one coproduct component.  The traced path family is
therefore the disjoint union of the path family of \(\Phi\) and the
untraced component \(\Psi\).  Hence
\[
    \operatorname{Tr}^{\mathbb C}_{\mathsf K}(\Phi)
    +
    \Psi
    =
    \operatorname{Tr}^{\mathbb C}_{\mathsf K}
    (\Phi+\Psi)
\]
under the canonical symmetry and associativity isomorphisms.

\textbf{Yanking.}
For the coproduct symmetry
\[
    \sigma_{\mathbb C,\mathbb C}:
    \mathbb C+\mathbb C
    \longrightarrow
    \mathbb C+\mathbb C,
\]
the direct visible block is initial, the entrance and exit blocks are
identities, and the hidden loop block is initial.  Only the zero-length
hidden path contributes, so
\[
    \operatorname{Tr}^{\mathbb C}_{\mathsf K}
    (\sigma_{\mathbb C,\mathbb C})
    =
    1_{\mathbb C}.
\]

All comparisons used above are pure invertible effectful simulations.
Passing to isomorphism classes makes the trace equations literal.
\end{proof}

\subsection{Trace on simulations}
\label{subsec:effectful-trace-simulations-uniformity}

Let
\[
    \widehat\Theta:
    \Phi\Rightarrow\Psi
\]
be a central effectful simulation between
\[
    \Phi,\Psi:
    \mathbb A+\mathbb C
    \rightsquigarrow_{\mathsf K}
    \mathbb B+\mathbb C.
\]
By Proposition~\ref{prop:effectful-block-restriction}, block restriction gives
simulations
\[
    \widehat\Theta_{AB},
    \qquad
    \widehat\Theta_{AC},
    \qquad
    \widehat\Theta_{CB},
    \qquad
    \widehat\Theta_{CC}.
\]

Define
\[
    \widehat\Theta_{CC}^{[0]}
    :=
    1_{1_{\mathbb C}},
\]
and recursively
\[
    \widehat\Theta_{CC}^{[n+1]}
    :=
    \widehat\Theta_{CC}^{[n]}
    \ast_h
    \widehat\Theta_{CC}.
\]

\begin{definition}[Trace of an effectful simulation]
\label{def:trace-effectful-simulation}
Define
\[
    \operatorname{Tr}^{\mathbb C}_{\mathsf K}
    (\widehat\Theta)
    :=
    \widehat\Theta_{AB}
    \boxplus
    \boxplus_{n\geq0}
    \left(
        \widehat\Theta_{CB}
        \ast_h
        \widehat\Theta_{CC}^{[n]}
        \ast_h
        \widehat\Theta_{AC}
    \right).
\]
\end{definition}

\begin{theorem}[Hom-category effectful trace]
\label{thm:hom-category-effectful-trace}
The assignments of
Definitions~\ref{def:effectful-coproduct-trace} and
\ref{def:trace-effectful-simulation} define functors on the hom-categories.
The trace coherence equations are witnessed by canonical invertible central
effectful simulations.
\end{theorem}

\begin{proof}
Block restriction, horizontal composition, and homwise sums preserve central
simulations by
Propositions~\ref{prop:effectful-block-restriction},
\ref{prop:horizontal-effectful-simulation-law}, and
\ref{prop:homwise-sums-effectful-mealy}.  Thus the displayed trace is a
central simulation.

For vertically composable simulations
\[
    \widehat\Theta:\Phi\Rightarrow\Psi,
    \qquad
    \widehat\Lambda:\Psi\Rightarrow\Xi,
\]
middle-four interchange gives, by induction,
\[
    (\widehat\Lambda\circ_v\widehat\Theta)_{CC}^{[n]}
    =
    \widehat\Lambda_{CC}^{[n]}
    \circ_v
    \widehat\Theta_{CC}^{[n]}.
\]
The same equation holds for the entrance, exit, and visible blocks.  Taking
the homwise coproduct over all hidden lengths and using the vertical
composition equation of
Proposition~\ref{prop:homwise-sums-effectful-mealy} gives
\[
    \operatorname{Tr}
    (\widehat\Lambda\circ_v\widehat\Theta)
    =
    \operatorname{Tr}(\widehat\Lambda)
    \circ_v
    \operatorname{Tr}(\widehat\Theta).
\]
Identity simulations are preserved similarly.

The coherence simulations are the pure images of the decorated hidden-path
reindexing isomorphisms from
Lemma~\ref{lem:decorated-hidden-path-normal-form}.
\end{proof}

\subsection{Trace on the retractive presentation quotient}
\label{subsec:trace-retractive-presentation-quotient}

\begin{proposition}[Coproduct stability of strict state retractions]
\label{prop:coproduct-stability-strict-state-retractions}
If
\[
    h:\Phi\to\Phi',
    \qquad
    k:\Psi\to\Psi'
\]
are strict state retractions, then
\[
    h+k:
    \Phi+\Psi
    \longrightarrow
    \Phi'+\Psi'
\]
is a strict state retraction.
\end{proposition}

\begin{proof}
Let \(s\) and \(t\) be strict sections of \(h\) and \(k\).  The strict
transition equation for \(h+k\) and for \(s+t\) can be checked after
restriction to the two coproduct summands.  Moreover,
\[
    (h+k)\circ(s+t)
    =
    (h\circ s)+(k\circ t)
    =
    1.
\]
\end{proof}

\begin{proposition}[Trace preserves strict state homomorphisms]
\label{prop:trace-preserves-strict-state-homomorphisms}
A strict state homomorphism
\[
    h:\Phi\longrightarrow\Psi
\]
induces a strict state homomorphism
\[
    \operatorname{Tr}^{\mathbb C}_{\mathsf K}(h):
    \operatorname{Tr}^{\mathbb C}_{\mathsf K}(\Phi)
    \longrightarrow
    \operatorname{Tr}^{\mathbb C}_{\mathsf K}(\Psi).
\]
\end{proposition}

\begin{proof}
By Proposition~\ref{prop:effectful-block-restriction}, restrict \(h\) to the
four blocks.  Form the iterated strict maps
\[
    h_{CC}^{[n]}:
    \Phi_{CC}^{[n]}
    \longrightarrow
    \Psi_{CC}^{[n]}.
\]
On the \(n\)-th hidden summand, use
\[
    h_{CB}
    \circ
    h_{CC}^{[n]}
    \circ
    h_{AC}.
\]
Together with \(h_{AB}\), these maps copair to
\[
    \operatorname{Tr}^{\mathbb C}_{\mathsf K}(h).
\]
The strict transition equation holds on every summand and therefore on their
coproduct by
Lemma~\ref{lem:strict-presentation-homwise-coproducts}.
\end{proof}

\begin{proposition}[Trace preserves strict state retractions]
\label{prop:trace-preserves-strict-state-retractions}
If
\[
    h:\Phi\longrightarrow\Psi
\]
is a strict state retraction, then
\[
    \operatorname{Tr}^{\mathbb C}_{\mathsf K}(h):
    \operatorname{Tr}^{\mathbb C}_{\mathsf K}(\Phi)
    \longrightarrow
    \operatorname{Tr}^{\mathbb C}_{\mathsf K}(\Psi)
\]
is a strict state retraction.
\end{proposition}

\begin{proof}
Let
\[
    s:\Psi\longrightarrow\Phi
\]
be a strict section, so that
\[
    h\circ s=1_\Psi.
\]
By Proposition~\ref{prop:effectful-block-restriction}, each \(s_{XY}\) is a
strict section of \(h_{XY}\).  By
Proposition~\ref{prop:trace-preserves-strict-state-homomorphisms},
\(\operatorname{Tr}(s)\) is strict.

On the visible summand,
\[
    h_{AB}\circ s_{AB}=1.
\]
On the \(n\)-th hidden-path summand, functoriality of pipeline composition
gives
\[
\begin{aligned}
&
\left(
    h_{CB}\circ h_{CC}^{[n]}\circ h_{AC}
\right)
\circ
\left(
    s_{CB}\circ s_{CC}^{[n]}\circ s_{AC}
\right)
\\
&\qquad=
1.
\end{aligned}
\]
Compatibility of homwise coproducts with vertical composition therefore
gives
\[
    \operatorname{Tr}(h)
    \circ
    \operatorname{Tr}(s)
    =
    1_{\operatorname{Tr}(\Psi)}.
\]
Thus \(\operatorname{Tr}(h)\) is a strict state retraction.
\end{proof}

\begin{theorem}[Traced retractive state-presentation quotient]
\label{thm:traced-retractive-state-presentation-quotient}
The category
\[
    \mathbf{Mly}^{\mathrm{cen},\mathrm{ret}}_{\mathsf K}(\mathcal E)
\]
inherits a coproduct symmetric monoidal structure and a finite-hidden-path
trace.

The canonical connected-component projection from
\[
    \mathbf{RetPres}^{\mathrm{cen}}_{\mathsf K}(\mathcal E)
\]
to the locally discrete bicategory determined by
\[
    \mathbf{Mly}^{\mathrm{cen},\mathrm{ret}}_{\mathsf K}(\mathcal E)
\]
preserves coproduct and trace.
\end{theorem}

\begin{proof}
Proposition~\ref{prop:coproduct-stability-strict-state-retractions} shows
that retractive presentation equivalence is a congruence for the coproduct
tensor.
Proposition~\ref{prop:trace-preserves-strict-state-retractions} shows that it
is a congruence for trace.

For every Joyal--Street--Verity trace equation, the decorated hidden-path
normal-form comparison of
Lemma~\ref{lem:decorated-hidden-path-normal-form} is a pure strict state
isomorphism and therefore a strict state retraction.  Hence the two sides of
each trace equation lie in the same connected component of the relevant
retractive-presentation hom-category.

The coproduct associator, unitors, and symmetry are likewise pure strict state
isomorphisms.  Therefore the symmetric monoidal and trace equations hold in
the retractive state-presentation quotient.
\end{proof}

\section{The one-object specialization}
\label{sec:one-object-effectful-mealy}

Let \(M\) and \(N\) be monoid objects in \(\mathcal E\), regarded as
one-object internal categories with the convention
\[
    b\circ a=ba.
\]
Since
\[
    \mathbf{Span}(\mathcal E)(1,1)
    \simeq
    \mathcal E/1
    \simeq
    \mathcal E,
\]
a carrier is an object \(P\), and an effectful Mealy morphism has an
underlying map
\[
    \widehat\Phi:
    M\times P
    \longrightarrow
    \mathsf K_1(P\times N).
\]

\begin{proposition}[One-object effectful Mealy morphisms]
\label{prop:one-object-effectful-mealy}
A one-object effectful Mealy morphism
\[
    M\rightsquigarrow_{\mathsf K}N
\]
is equivalently an object \(P\) and a map
\[
    \widehat\Phi:
    M\times P
    \to
    \mathsf K_1(P\times N)
\]
satisfying
\[
    \widehat\Phi(1_M,x)
    =
    \eta^{\mathsf K}(x,1_N)
\]
and
\[
\begin{aligned}
    \widehat\Phi(mn,x)
    ={}&
    \operatorname{bind}_{\mathsf K}
    \Bigl(
        \widehat\Phi(m,x),
        (x_1,u)
        \longmapsto
    \\
    &\qquad
        \operatorname{bind}_{\mathsf K}
        \Bigl(
            \widehat\Phi(n,x_1),
            (x_2,v)
            \longmapsto
            \eta^{\mathsf K}
            (x_2,uv)
        \Bigr)
    \Bigr).
\end{aligned}
\]
\end{proposition}

\begin{proof}
In the one-object case,
\[
    R_M(P)=M\times P,
    \qquad
    T_N(P)=P\times N.
\]
Apply Proposition~\ref{prop:effectful-mealy-component-laws} and the standing
one-object multiplication convention.
\end{proof}

A central one-object machine is one for which
\[
    M\times P
    \longrightarrow
    \mathsf K_1(P\times N)
\]
is fibred-central.  If every slice monad
\[
    \mathsf K_Z
\]
is commutative, then every one-object effectful Mealy morphism is central.
Commutativity of \(\mathsf K_1\) alone implies only fibrewise centrality in
the terminal slice.

An effectful simulation \(P\Rightarrow Q\) is a map
\[
    Q\longrightarrow\mathsf K_1(P\times N)
\]
satisfying the one-object specialization of
Proposition~\ref{prop:effectful-simulation-component-equation}.  Pipeline
composition is the bind formula of
Definition~\ref{def:effectful-pipeline-transition}.

\section{The stateless specialization}
\label{sec:stateless-effectful-mealy}

Let
\[
    F_0:A_0\longrightarrow B_0
\]
be a morphism, regarded as the map-representable carrier span
\[
    A_0
    \xleftarrow{1_{A_0}}
    A_0
    \xrightarrow{F_0}
    B_0.
\]
Then
\[
    P\circ A\cong A_1
\]
as the object
\[
    \left(
        A_1
        \xrightarrow{\langle s_A,F_0t_A\rangle}
        A_0\times B_0
    \right),
\]
while
\[
    B\circ P
    \cong
    A_0\times_{F_0,B_0,s_B}B_1
\]
as the object
\[
    \left(
        A_0\times_{F_0,B_0,s_B}B_1
        \xrightarrow{\langle\pi_1,t_B\pi_2\rangle}
        A_0\times B_0
    \right).
\]

\begin{proposition}[Stateless effectful Mealy morphisms]
\label{prop:stateless-effectful-mealy}
A stateless effectful Mealy morphism with object map \(F_0\) is equivalently
a map in \(\mathcal E/(A_0\times B_0)\)
\[
\widehat F_1:
\left(
    A_1,
    \langle s_A,F_0t_A\rangle
\right)
\longrightarrow
\mathsf K_{A_0\times B_0}
\left(
    A_0\times_{F_0,B_0,s_B}B_1,
    \langle\pi_1,t_B\pi_2\rangle
\right)
\]
satisfying the effectful identity and composition equations.
\end{proposition}

\begin{proof}
Substitute the map-representable carrier into
Definition~\ref{def:effectful-mealy-morphism}.  The updated state is forced
to be \(s_Aa\).  The remaining output is an effectful \(B\)-arrow
\[
    F_0(s_Aa)\longrightarrow F_0(t_Aa).
\]
The Mealy algebra equations become the effectful preservation equations for
identities and composition.
\end{proof}

Under the pure image, an internal functor
\[
    F:\mathbb A\longrightarrow\mathbb B
\]
produces
\[
    \widehat F_1(a)
    =
    \eta^{\mathsf K}
    \bigl(
        s_Aa,
        F_1a
    \bigr).
\]

Every internal natural transformation
\[
    \alpha:F\Rightarrow G
\]
determines, by application of \(\eta^{\mathsf K}\), a pure effectful
simulation between the corresponding pure stateless effectful Mealy
morphisms.  The converse is not asserted: the monad unit need not reflect
equality of ordinary output-comparison cells.

\section{Examples of internal effects}
\label{sec:examples-internal-effects}

\subsection{The identity effect}
\label{subsec:identity-effect-example}

Take
\[
    \mathsf K_Z
    =
    \operatorname{Id}_{\mathcal E/Z}
\]
for every \(Z\).  The two Fubini maps are identities, every computation is
central, and
\[
    \mathbf{Mly}^{\mathrm{cen}}_{\mathsf K}(\mathcal E)
    \cong
    \mathbf{Mly}(\mathcal E).
\]
The component laws reduce to the deterministic transition and output
equations.

\subsection{Reader effects}
\label{subsec:reader-effects}

Assume that \(\mathcal E\) is locally cartesian closed, and fix
\(R\in\mathcal E\).  For every \(Z\), put
\[
    R_Z
    :=
    (Z\times R\xrightarrow{\pi_Z}Z).
\]

\begin{definition}[Fibred reader monad]
\label{def:fibred-reader-monad}
Define
\[
    \mathsf{Reader}_{R,Z}(P)
    :=
    P^{R_Z}
\]
in the cartesian closed slice \(\mathcal E/Z\).
\end{definition}

The unit is the transpose of
\[
    P\times_ZR_Z\to P,
\]
and multiplication is internally
\[
    \mu(F)(r)=F(r)(r).
\]
Pullback reindexing preserves the slice exponentials and their evaluation
maps, so these monads assemble into a fibred strong monad.

\begin{proposition}[Centrality of reader effects]
\label{prop:reader-effects-central}
The fibred reader monad is commutative.  Hence every reader-effectful Mealy
morphism and simulation is central.
\end{proposition}

\begin{proof}
Internally in \(\mathcal E/Z\), both Fubini maps send
\[
    (f,g):
    P^{R_Z}\times_ZQ^{R_Z}
\]
to
\[
    r\longmapsto(f(r),g(r)).
\]
Thus
\[
    d^\ell=d^r.
\]
\end{proof}

A reader-effectful transition has the internal form
\[
    \widehat\Phi(a,x):
    R\longrightarrow
    \operatorname{Out}_P(s_Aa,qx).
\]
Every stage of sequential execution is evaluated at the same reader value.

\subsection{Writer effects}
\label{subsec:writer-effects}

Let
\[
    \mathbb W=(W,e_W,m_W)
\]
be a monoid object of \(\mathcal E\).  For each \(Z\), put
\[
    W_Z:=(Z\times W\to Z).
\]

\begin{definition}[Fibred writer monad]
\label{def:fibred-writer-monad}
Define
\[
    \mathsf{Writer}_{W,Z}(P)
    :=
    W_Z\times_ZP.
\]
The unit is
\[
    p\longmapsto(e_W,p),
\]
and multiplication is
\[
    (w,(v,p))
    \longmapsto
    (m_W(w,v),p).
\]
\end{definition}

Pullback preserves these products, so the writer monads form a fibred strong
monad.

Assume that \(\mathcal E\) is cartesian closed.  Let
\[
    \lambda,\rho:W\longrightarrow W^W
\]
be the transposes of left and right multiplication:
\[
    \lambda(w)(v)=m_W(w,v),
    \qquad
    \rho(w)(v)=m_W(v,w).
\]
Define the internal centre by the equalizer
\[
\begin{tikzcd}
    Z(W)
        \arrow[r,hook,"z_W"]
    &
    W
        \arrow[r,shift left=.7ex,"\lambda"]
        \arrow[r,shift right=.7ex,"\rho"']
    &
    W^W.
\end{tikzcd}
\]

\begin{proposition}[Central writer computations]
\label{prop:central-writer-computations}
A writer Kleisli arrow
\[
    f:P\longrightarrow W_Z\times_ZQ
\]
is fibrewise central if and only if its writer component
\[
    w_f:P\longrightarrow W
\]
factors through
\[
    z_W:Z(W)\hookrightarrow W.
\]
Such a factorization is stable under reindexing, so fibrewise and fibred
centrality coincide for writer arrows.

The full writer monad is commutative if and only if \(W\) is a commutative
monoid object.
\end{proposition}

\begin{proof}
Let
\[
    g:R\longrightarrow W_Z\times_ZS
\]
have writer component \(w_g\).  The two Fubini maps produce writer values
\[
    m_W(w_f,w_g)
    \quad\text{and}\quad
    m_W(w_g,w_f).
\]
Equality for every \(g\) is equivalent, internally, to
\[
    m_W(w_f(p),v)
    =
    m_W(v,w_f(p))
\]
for every \(v\in W\).  This is precisely the equality
\[
    \lambda w_f=\rho w_f,
\]
and therefore precisely factorization through the equalizer \(Z(W)\).

Pullback preserves the equalizer and the displayed factorization, so the
criterion is reindexing-stable.  The monad is commutative exactly when the
identity \(W\to W\) factors through \(Z(W)\), equivalently when
\(\lambda=\rho\).
\end{proof}

A central writer-effectful transition produces
\[
    (w,x',\beta)
\]
with \(w\) in the internal centre of \(W\).  Sequential execution multiplies
writer values in execution order and composes output arrows.

\subsection{Partiality}
\label{subsec:partiality-effect}

Assume that \(\mathcal E\) has pullback-stable finite coproducts.  The
terminal object of \(\mathcal E/Z\) is
\[
    1_Z=(Z\xrightarrow{1_Z}Z).
\]

\begin{definition}[Fibred lift monad]
\label{def:fibred-lift-monad}
Define
\[
    \mathsf L_Z(P)
    :=
    P+_Z1_Z
\]
in \(\mathcal E/Z\).  The unit is the left coproduct injection, and
multiplication
\[
    (P+_Z1_Z)+_Z1_Z
    \longrightarrow
    P+_Z1_Z
\]
is the identity on \(P\) and the codiagonal on the two failure summands.
\end{definition}

Since coproducts in the slice are computed by coproducts in \(\mathcal E\),
one has concretely
\[
    \mathsf L_Z(P)
    =
    \bigl(
        P+Z\xrightarrow{[p,1_Z]}Z
    \bigr).
\]
When the coproducts are disjoint, this object decomposes into a success
summand and a failure summand in the expected way.  Failure retains the
external interface \(Z\).

\begin{proposition}[Centrality of partiality]
\label{prop:partiality-central}
The fibred lift monad is commutative.  Hence every partial effectful Mealy
morphism and simulation is central.
\end{proposition}

\begin{proof}
The product of two lifted computations decomposes into four summands:
success--success, success--failure, failure--success, and failure--failure.
On the success--success summand, both Fubini maps return the same ordered
pair.  On every other summand, both return the failure value over the current
interface.  Hence
\[
    d^\ell=d^r.
\]
\end{proof}

A partial transition either returns a typed pair \((x',\beta)\) or the
failure point over the same external interface.  Identity inputs succeed, and
a composite succeeds exactly when both constituent transitions succeed.

\subsection{Exception effects}
\label{subsec:exception-effects}

Let \(E\in\mathcal E\), and assume pullback-stable finite coproducts.  Put
\[
    E_Z:=(Z\times E\to Z)
\]
and define
\[
    \mathsf{Exc}_{E,Z}(P)
    :=
    P+_ZE_Z.
\]
The unit inserts a successful value.  Multiplication keeps the first
exception encountered in execution order.

\begin{proposition}[Commutativity criterion for exceptions]
\label{prop:exception-commutativity}
If
\[
    E\longrightarrow1
\]
is subterminal, then the fibred exception monad is commutative.

Conversely, commutativity implies
\[
    \iota_E\pi_1
    =
    \iota_E\pi_2:
    E\times E\longrightarrow1+E,
\]
where \(\iota_E:E\to1+E\) is the exception injection.  Consequently, when
the relevant coproduct injections are monic, and in particular in a feedback
span base, the fibred exception monad is commutative if and only if
\(E\to1\) is subterminal.
\end{proposition}

\begin{proof}
If \(E\) is subterminal, any two exception values in the same context agree.
Therefore the order in which failing computations are executed is
irrelevant, and the two Fubini maps agree.

Conversely, apply commutativity to the two generic exception computations
over \(E\times E\).  One Fubini map returns the first exception and the other
returns the second.  Equality gives
\[
    \iota_E\pi_1
    =
    \iota_E\pi_2:
    E\times E\to1+E.
\]
If \(\iota_E\) is monic, this implies
\[
    \pi_1=\pi_2:
    E\times E\to E,
\]
which is equivalent to subterminality of \(E\to1\).  In a feedback span base,
disjointness of coproducts makes the coproduct injections monic.
\end{proof}

For a general exception object, the central exception computations are
exactly those satisfying
Definition~\ref{def:fibred-central-kleisli-arrow}.

\subsection{State effects}
\label{subsec:state-effects}

Assume that \(\mathcal E\) is locally cartesian closed, and fix a store object
\(S\).  Put
\[
    S_Z:=(Z\times S\to Z).
\]

\begin{definition}[Fibred state monad]
\label{def:fibred-state-monad}
Define
\[
    \mathsf{State}_{S,Z}(P)
    :=
    (S_Z\times_ZP)^{S_Z}
\]
in \(\mathcal E/Z\).
\end{definition}

Internally, an element is a state transformer
\[
    S\longrightarrow S\times P.
\]
The unit returns its value without changing the store, and multiplication
sequences the outer and inner state transformations.  Reindexing preserves
the products and slice exponentials, so these monads form a fibred strong
monad.

The two Fubini maps are the two execution orders for store transformations.
Thus a state-effectful transition is central precisely when it commutes with
every store computation after every reindexing.

\begin{proposition}[Central state computations in \(\mathbf{Set}\)]
\label{prop:central-state-computations-set}
Let \(\mathcal E=\mathbf{Set}\), and let \(S\) be nonempty.  A Kleisli map
\[
    f:P\longrightarrow(S\times Q)^S
\]
is central if and only if it is pure.  Equivalently, there is a unique map
\[
    \overline f:P\longrightarrow Q
\]
such that
\[
    f(p)(s)
    =
    (s,\overline f(p))
\]
for every \(p\in P\) and \(s\in S\).
\end{proposition}

\begin{proof}
Write
\[
    f(p)(s)
    =
    (u_p(s),x_p(s)).
\]
Let an arbitrary comparison computation be
\[
    g(r)(s)
    =
    (v_r(s),y_r(s)).
\]
Centrality gives
\[
    v_r(u_p(s))
    =
    u_p(v_r(s)),
\]
\[
    x_p(s)
    =
    x_p(v_r(s)),
\]
and
\[
    y_r(u_p(s))
    =
    y_r(s).
\]

First take a comparison computation with domain \(1\), returned value in
\(S\), and
\[
    g(*)(s)=(s,s).
\]
Its returned-value component is \(1_S\), so the third equation gives
\[
    u_p(s)=s.
\]

Now fix \(t\in S\), and choose a computation whose store-update component is
the constant map
\[
    v(s)=t.
\]
The second equation gives
\[
    x_p(s)=x_p(t).
\]
Since \(s,t\) are arbitrary, \(x_p\) is constant.  Define its value to be
\(\overline f(p)\).  This proves that \(f\) is pure.  The converse follows
because pure Kleisli arrows are central.
\end{proof}

\subsection{Set-based branching and quantitative effects}
\label{subsec:set-based-branching-effects}

Now specialize to
\[
    \mathcal E=\mathbf{Set}.
\]
For a Set-monad \(K\), define the associated fibred monad on a family
\(P\to Z\) by
\[
    (\mathsf K_ZP)_z
    =
    K(P_z).
\]
Equivalently,
\[
    \mathsf K_Z(P)
    =
    \coprod_{z\in Z}K(P_z)
    \longrightarrow
    Z.
\]

\begin{proposition}[Fibrewise Set construction]
\label{prop:fibrewise-set-construction}
Every strong monad \(K\) on \(\mathbf{Set}\) induces a fibred strong monad on
the codomain fibration by fibrewise application.  If \(K\) is commutative,
every effectful Mealy morphism and simulation for this fibred monad is
central.
\end{proposition}

\begin{proof}
Pullback of a family is reindexing of its fibres.  Fibrewise application of
\(K\) commutes strictly with this reindexing.  The monad operations and
strength are defined on each fibre.  If the two Fubini maps of \(K\) agree,
they agree in every fibre and therefore in every slice.
\end{proof}

\subsubsection{Nondeterminism}

For the finite-powerset monad \(\mathcal P_{\mathrm f}\), an effectful
transition is a finite set of typed outcomes:
\[
    \widehat\Phi(a,x)
    \subseteq
    \operatorname{Out}_P(s_Aa,qx).
\]
The component laws are
\[
    \widehat\Phi(e_A(px),x)
    =
    \{(x,e_B(qx))\}
\]
and
\[
\begin{aligned}
\widehat\Phi(b\circ a,x)
=
\bigl\{
    (x_a,\beta_b\circ\beta_a)
    \,\big|\,
    &(x_b,\beta_b)\in\widehat\Phi(b,x),
    \\
    &(x_a,\beta_a)\in\widehat\Phi(a,x_b)
\bigr\}.
\end{aligned}
\]
The finite-powerset and full-powerset monads are commutative, so all such
machines are central.

\subsubsection{Probability and subprobability}

For the finite-distribution monad \(\mathcal D_{\mathrm f}\),
\(\widehat\Phi(a,x)\) is a finitely supported probability distribution on
typed pairs \((x',\beta)\).  The identity input has the Dirac distribution at
\[
    (x,e_B(qx)).
\]
For composable \(a,b\),
\[
\begin{aligned}
&
\widehat\Phi(b\circ a,x)(x_a,\gamma)
\\
&\quad=
\sum_{\substack{
    x_b,\beta_b,\beta_a\\
    \beta_b\circ\beta_a=\gamma
}}
\widehat\Phi(b,x)(x_b,\beta_b)
\widehat\Phi(a,x_b)(x_a,\beta_a).
\end{aligned}
\]
This is the typed Chapman--Kolmogorov convolution.

The finite-distribution and finite-subdistribution monads are commutative.
The subdistribution monad permits loss of mass and therefore models failure
or nontermination.

\subsubsection{Semiring-weighted effects}

Let \(R\) be a semiring, and let \(\mathsf W_R(X)\) be the set of finitely
supported functions \(X\to R\).  If
\[
    w_{a,x}(x',\beta)\in R
\]
is the weight of an outcome, then
\[
\begin{aligned}
w_{b\circ a,x}(x_a,\gamma)
=
\sum_{\substack{
    x_b,\beta_b,\beta_a\\
    \beta_b\circ\beta_a=\gamma
}}
w_{b,x}(x_b,\beta_b)
w_{a,x_b}(x_a,\beta_a).
\end{aligned}
\]

\begin{proposition}[Central weighted computations]
\label{prop:central-weighted-computations}
A finitely supported \(R\)-weighted computation is central if and only if all
of its coefficients lie in the multiplicative centre
\[
    Z(R)
    =
    \{r\in R\mid rs=sr\text{ for every }s\in R\}.
\]
Hence the weighted monad is commutative exactly when multiplication in
\(R\) is commutative.
\end{proposition}

\begin{proof}
For outcomes with coefficients \(r\) and \(s\), the two Fubini maps assign
coefficients
\[
    rs
    \quad\text{and}\quad
    sr.
\]
If every coefficient of the first computation is central, these agree for
every second computation.

Conversely, compare with a singleton-supported computation of arbitrary
coefficient \(s\).  Equality at each ordered pair of support points forces
every coefficient \(r\) of the first computation to satisfy
\[
    rs=sr.
\]
\end{proof}

The Boolean semiring recovers finite nondeterminism, the natural-number
semiring counts executions with multiplicity, and commutative coefficient
semirings yield fully central weighted Mealy bicategories.

\section{Chapter summary}
\label{sec:effectful-chapter-summary}

This chapter constructed effectful Mealy morphisms from a fibred strong monad
on the codomain fibration of a finitely complete category.

\begin{enumerate}
    \item Every interface object \(Z\) carries a slice monad
    \[
        \mathsf K_Z,
    \]
    and every map \(h:W\to Z\) induces a dependent-sum Kleisli functor
    \[
        \Sigma_h^{\mathsf K}:
        \operatorname{Kl}(\mathsf K_W)
        \to
        \operatorname{Kl}(\mathsf K_Z).
    \]

    \item The dependent-sum transport satisfies functoriality,
    Beck--Chevalley, compatibility with left and right strength, and the
    effectful projection formula.

    \item The strength gives two Fubini maps
    \[
        d^\ell,d^r:
        \mathsf KP\times\mathsf KQ
        \to
        \mathsf K(P\times Q).
    \]
    Their discrepancy is exactly the obstruction to unrestricted
    middle-four interchange.

    \item Pure parameters define genuine left and right Kleisli functors.
    After reindexing and dependent-sum transport, these give functorial pure
    whiskering of effectful span cells and the coherence needed for the
    canonical distributive laws.

    \item Fibrewise central arrows commute with every computation in one
    slice.  Fibred-central arrows are those whose every reindexing is
    fibrewise central.  They are closed under Kleisli composition, central
    tensor, reindexing, and dependent-sum transport.

    \item Objects of \(\mathcal E\), spans, and central effectful span cells
    form the bicategory
    \[
        \mathbf{Span}^{\mathrm{cen}}_{\mathsf K}(\mathcal E).
    \]

    \item Unrestricted effectful span cells form a bicategory if and only if
    every slice monad is commutative.  The converse follows from a direct
    middle-four calculation using the two sequential decompositions of the
    Fubini maps.

    \item Pure whiskering constructs canonical distributive laws
    \[
        T_B\mathsf K\Rightarrow\mathsf KT_B,
        \qquad
        R_A\mathsf K\Rightarrow\mathsf KR_A.
    \]
    Together with the input--output distributive law of Chapter~2, they
    satisfy Yang--Baxter compatibility by mixed associativity of pure
    whiskering.

    \item An effectful Mealy morphism is an algebra for the lifted input
    monad, with structure
    \[
        P\circ A
        \longrightarrow
        \mathsf K(B\circ P).
    \]
    A central such structure is equivalently an oriented monad morphism in
    the central effectful span bicategory.

    \item Central effectful Mealy morphisms and central effectful simulations
    form the bicategory
    \[
        \mathbf{Mly}^{\mathrm{cen}}_{\mathsf K}(\mathcal E).
    \]

    \item Strict state homomorphisms form a covariant presentation
    bicategory.  Under a pullback-stable initial object, the connected
    component quotient generated by all strict state homomorphisms collapses
    every hom-category.

    \item Strict state retractions form a locally wide presentation
    bicategory.  In a feedback span base, the empty retractive class is
    isolated from every noninitial carrier, so the empty presentation no
    longer forces the quotient to collapse.

    \item The connected-component quotient
    \[
        \mathbf{Mly}^{\mathrm{cen},\mathrm{ret}}_{\mathsf K}(\mathcal E)
    \]
    is universal for typed retract-invariant unpointed state hiding.

    \item The retractive state-hiding quotient is distinct from pointed
    uniform feedback: an initial state is not part of the present Mealy
    data.

    \item When \(\mathcal E\) is a feedback span base, homwise sums of
    effectful Mealy morphisms and simulations exist, preserve centrality, and
    distribute over pipeline composition.

    \item Coproduct lifts to a symmetric monoidal structure on central
    effectful Mealy morphisms.

    \item For a central effectful endomorphism \(U\),
    \[
        \operatorname{Path}_{\mathsf K}(U)
        =
        \boxplus_{n\geq0}U^{[n]}
    \]
    is the free path monoid in the category of strict state presentations.

    \item The finite hidden-path trace is
    \[
        \operatorname{Tr}^{\mathbb C}_{\mathsf K}(\Phi)
        =
        \Phi_{AB}
        \boxplus
        \Phi_{CB}
        \circ
        \operatorname{Path}_{\mathsf K}(\Phi_{CC})
        \circ
        \Phi_{AC}.
    \]
    It acts functorially on simulations, preserves strict state
    homomorphisms and strict state retractions, and descends to the retractive
    state-presentation quotient.

    \item Reader, writer, partiality, exception, state, nondeterministic,
    probabilistic, subprobabilistic, and semiring-weighted effects arise
    internally from the same fibred construction.
\end{enumerate}

The structural relation established in the chapter is
\[
    \boxed{
    \begin{gathered}
        \text{fibred strong effects}
        \quad+\quad
        \text{reindexing-stable centrality}
        \quad+\quad
        \text{span-carried state}
        \\
        =
        \text{composable effectful Mealy systems with finite-path feedback}.
    \end{gathered}
    }
\]